\subsection{结点向量与逆根}

引理 2. 设 S 是 T 的闭子群，Z(S) 是它在 G 中的正规化子。

\begin{enumerate}
\def\labelenumi{(\roman{enumi})}
\item
  \(\mathrm{R}(\mathrm{Z}(\mathrm{S})_0,\mathrm{T})\) 是诸 \(\alpha \in \mathrm{R}(\mathrm{G},\mathrm{T})\) 中满足 \(\alpha (\mathrm{S}) = \{1\}\) 者所成的集合。
\item
  \(\mathrm{Z(S)_0}\) 的中心是诸 \(\operatorname {Ker}\alpha\)（\(\alpha \in \mathrm{R}(\mathrm{Z(S)}_0,\mathrm{T})\)）的交。
\item
  若 S 连通，则 Z(S) 连通。
\end{enumerate}

李代数 \(\mathrm{L}(\mathrm{Z}(\mathrm{S}))_{(\mathbf{C})}\) 由 S 在 \(g_{C}\) 上的不变量组成（第三章，§9，第 3 节，命题 8），因而是 \(t_{C}\) 与诸 \(g^{\alpha}\)（满足 \(\alpha(\mathrm{S})=\{1\}\) 者）的直和，由此得 (i)。断言 (ii) 由命题 8（第 4 节）得出，而断言 (iii) 已在 §2 第 2 节定理 2 的推论 5 中证明。

定理 1. 设 \(\alpha \in \mathrm{R}(\mathrm{G},\mathrm{T})\)。令 \(Z_{\alpha}\) 为 \(\alpha\) 的核的中心化子；它是 \(\mathrm{G}\) 的连通闭子群；其中心是 \(\operatorname{Ker} \alpha\)；其换位子群 \(\mathrm{D}(Z_{\alpha}) = S_{\alpha}\) 是 \(\mathrm{G}\) 的秩 1 连通闭半单子群。我们有 \(\mathrm{R}(Z_{\alpha},\mathrm{T}) = \{\alpha, -\alpha\}\)，且 \(\dim Z_{\alpha} = \dim T + 2\)。

设 \(Z_{\alpha}^{\prime}\) 是 \((\mathrm{Ker}\alpha)_0\) 的中心化子。由引理 2，它是 G 的连通闭子群，且 \(\mathrm{R}(Z_{\alpha}^{\prime},\mathrm{T})\) 是 \(\beta \in \mathrm{R}(\mathrm{G},\mathrm{T})\) 中满足 \(\beta ((\mathrm{Ker}\alpha)_0) = \{1\}\) 者所成的集合。显然，\(\{\alpha , - \alpha \} \subset \mathrm{R}(Z_{\alpha}^{\prime},\mathrm{T})\)。反之，设 \(\beta \in \mathrm{R}(Z_{\alpha}^{\prime},\mathrm{T})\)；由于 \((\mathrm{Ker}\alpha)_0\) 在 \(\mathrm{Ker}\alpha\) 中指数有限，故存在整数 \(r\neq 0\)，使得 \(t^{r\beta} = 1\)（对 \(t\in \mathrm{Ker}\alpha\)）。由序列

\[
0 \longrightarrow \mathbf {Z} \longrightarrow \mathrm{X(T)} \longrightarrow \mathrm{X(Ker} \alpha) \longrightarrow 0
\]

（它通过对偶对应于正合序列

\[
0 \longrightarrow \operatorname{Ker} \alpha \longrightarrow \mathrm{T} \stackrel {{\alpha}} {{\longrightarrow}} \mathrm{U} \longrightarrow 0,
\]

）的正合性可知，\(r\beta\) 是 \(\alpha\) 的倍数；由第八章 §2 第 2 节定理 2 (i)，这蕴含 \(\beta \in \{\alpha, -\alpha\}\)。于是 \(\mathrm{R}(Z_{\alpha}', T) = \{\alpha, -\alpha\}\)。由此并根据引理 2，\(Z_{\alpha}'\) 的中心是 \(\operatorname{Ker} \alpha\)，故 \(Z_{\alpha}' = Z_{\alpha}\)。最后，由 §1 第 4 节命题 4 的推论 1，\(\mathrm{D}(\mathrm{Z}_{\alpha})\) 是 \(\mathrm{G}\) 的连通闭半单子群；它秩为 1，因为 \(\mathcal{DL}(\mathrm{Z}_{\alpha})_{(\mathbf{C})} = \mathfrak{g}^{\alpha} + \mathfrak{g}^{-\alpha} + [\mathfrak{g}^{\alpha},\mathfrak{g}^{-\alpha}]\)。

推论. 存在具有以下性质的李群态射 \(\nu : \mathbf{SU}(2, \mathbf{C}) \to \mathrm{G}\)：

\begin{enumerate}
\def\labelenumi{\alph{enumi})}
\item
  \(\nu\) 的像与 \(\alpha\) 的核交换。
\item
  对所有 \(a \in U\)，我们有 \(\nu\left(\begin{array}{cc}a & 0 \\ 0 & \bar{a}\end{array}\right) \in T\)，且 \(\alpha \circ \nu\left(\begin{array}{cc}a & 0 \\ 0 & \bar{a}\end{array}\right) = a^{2}\)。
\end{enumerate}

若 \(\nu_{1}\) 与 \(\nu_{2}\) 是从 \(\mathbf{SU}(2,\mathbf{C})\) 到 \(\mathrm{G}\) 且具有前述性质的两个态射，则存在 \(a\in \mathbf{U}\) 使得 \(\nu_{2} = \nu_{1}\circ \operatorname {Int}\left( \begin{array}{cc}a & 0\\ 0 & \bar{a} \end{array} \right)\)。

由定理 1 以及 §3 第 6 节命题 6，存在李群态射 \(\nu : \mathbf{SU}(2, \mathbf{C}) \to \mathrm{S}_{\alpha}\)，它是满射且核为离散。于是 \(\nu^{-1}(\mathrm{T} \cap \mathrm{S}_{\alpha})\) 是 \(\mathbf{SU}(2, \mathbf{C})\) 的一个极大环面（§2，第 3 节，命题 1）。由于 \(\mathbf{SU}(2, \mathbf{C})\) 的诸极大环面彼此共轭（§2，第 2 节，定理 2），必要时把 \(\nu\) 换为 \(\nu \circ \operatorname{Int}s\)（其中 \(s \in \mathbf{SU}(2, \mathbf{C})\)），可以假设 \(\nu^{-1}(\mathrm{T} \cap \mathrm{S}_{\alpha})\) 是 \(\mathbf{SU}(2, \mathbf{C})\) 中的对角矩阵群。于是 \(\nu\begin{pmatrix} a & 0 \\ 0 & \bar{a} \end{pmatrix} \in \mathrm{T}\) 对所有 \(a \in \mathbf{U}\) 成立，且映射

\[
\left( \begin{array}{c c} a & 0 \\ 0 & \bar {a} \end{array} \right) \mapsto \alpha \circ \nu \left( \begin{array}{c c} a & 0 \\ 0 & \bar {a} \end{array} \right)
\]

是 \(\mathbf{SU}(2,\mathbf{C})\) 的一个根，从而等于两个映射 \(\begin{pmatrix}a&0\\ 0&\bar{a}\end{pmatrix}\mapsto a^{2}\) 或 \(\begin{pmatrix}a&0\\ 0&\bar{a}\end{pmatrix}\mapsto a^{-2}\) 之一（§3，第 6 节，公式 (19)）。在第一种情形，同态 \(\nu\) 具有所要求的性质；在第二种情形，同态 \(\nu\circ Int\theta\) 具有这些性质（同前处，公式 (18)）。

若 \(\nu_{1}\) 与 \(\nu_{2}\) 是从 \(\mathbf{SU}(2,\mathbf{C})\) 到 G 且满足所述条件的两个态射，则它们都把 \(\mathbf{SU}(2,\mathbf{C})\) 映到 \(S_{\alpha}\) 内（条件 a)），因而都是 \(S_{\alpha}\) 的泛覆盖。于是，存在 \(\varphi\)（\(\mathbf{SU}(2,\mathbf{C})\) 的一个自同构）使得 \(\nu_{2} = \nu_{1} \circ \varphi\)，再应用第 4 节命题 9 即得结论。

由前述推论可知，同态 \(\nu_{\mathrm{T}}\)（从 \(\mathbf{U}\) 到 \(\mathrm{T}\)，由 \(\nu_{\mathrm{T}}(a) = \nu \left( \begin{array}{cc}a & 0\\ 0 & \bar{a} \end{array} \right)\)（\(a\in \mathbf{U}\)）定义）与 \(\nu\) 的选取无关。用 \(K_{\alpha}\in \Gamma (\mathrm{T})\) 表示在 \(\Gamma (\nu_{\mathrm{T}})\) 之下元素 \(2\pi i\)（视为 \(\Gamma (\mathbf{U}) = 2\pi i\mathbf{Z}\) 的元素）所对应的像；它称为相伴于根 \(\alpha\) 的结点向量。我们有 \(\langle \alpha ,K_{\alpha}\rangle = 2\)，换言之（第 2 节，公式 (2)）\(\delta (\alpha)(K_{\alpha}) = 4\pi i\)；由于 \(K_{\alpha}\) 属于 \(\mathfrak{t}\) 与 \(\mathrm{L}(\mathrm{S}_{\alpha})_{(\mathbf{C})}\) 的交，我们有

\[
K _ {\alpha} = 2 \pi i H _ {\delta (\alpha)},\tag{13}
\]

其中 \(H_{\delta(\alpha)}\) 是根 \(\delta(\alpha)\)（\((\mathfrak{g}_{\mathbf{C}}, \mathfrak{t}_{\mathbf{C}})\) 的根）所相伴的逆根（第八章，§2，第 2 节）。换言之，当 \(\Gamma(\mathrm{T}) \otimes \mathbf{R}\) 等同于 \(\mathrm{X(T)}\otimes \mathbf{R}\) 的对偶（借助配对 \(\langle ,\rangle ,K_{\alpha}\)）时，该元素便等同于逆根 \(\alpha^{\vee}\in (\mathrm{X(T)}\otimes \mathbf{R})^{*}\)。

注. 对所有 \(x \in R\)，我们有

\[
\nu \left( \begin{array}{c c} \exp (2 \pi i x) & 0 \\ 0 & \exp (- 2 \pi i x) \end{array} \right) = \nu_ {\mathrm{T}} (e ^ {2 \pi i x}) = \exp (x K _ {\alpha}).\tag{14}
\]

特别地：

\[
\nu \left( \begin{array}{c c} - 1 & 0 \\ 0 & - 1 \end{array} \right) = \nu_ {\mathrm{T}} (- 1) = \exp \Bigl (\frac {1}{2} K _ {\alpha} \Bigr).\tag{15}
\]

由此可知，\(\nu\) 是单射当且仅当 \(K_{\alpha} \notin 2\Gamma(T)\)，换言之，当且仅当存在 \(\lambda \in \mathrm{X}(T)\) 使得 \(\langle \lambda, K_{\alpha} \rangle \notin 2Z\)。当 \(g_{C}\) 为单时，\(\nu\) 都是单射，除非 \(g_{C}\) 是 \(B_{n}\) 型、\(C(G) = \{1\}\) 且 \(\alpha\) 是短根（参看第六章，图版）。

在本段余下部分，我们用 \(\mathrm{R}^{\vee}(\mathrm{G},\mathrm{T})\) 表示诸结点向量 \(K_{\alpha}\)（\(\alpha\in\mathrm{R}(\mathrm{G},\mathrm{T})\)）所成的集合。它是 \(\varGamma(\mathrm{T})\) 的一个子集，\(\varGamma(\mathrm{T})\) 到 \(t_{C}\) 中的典则单射把它等同于对比率为 \(2\pi i\) 的位似作用于逆根系 \(\mathrm{R}^{\vee}(\mathfrak{g}_{\mathbf{C}},\mathfrak{t}_{\mathbf{C}})=\{H_{\delta(\alpha)}\}\)（即 \(\delta(\mathrm{R})\) 的逆根系）所得的集合。由此可知，\(\mathrm{R}^{\vee}(\mathrm{G},\mathrm{T})\) 生成 R-向量空间 \(\mathrm{L}(\mathrm{T}\cap\mathrm{D}(\mathrm{G}))\)，从而它在 \(\mathrm{X}(\mathrm{T})\) 中的正交补是 \(\mathrm{X}(\mathrm{T}/(\mathrm{T}\cap\mathrm{D}(\mathrm{G})))\)。

用 \(\mathrm{Aut(T)}\) 表示李群 T 的自同构群；Weyl 群 \(W = W_{G}(T)\)（§2，第 5 节）可等同于 \(\mathrm{Aut(T)}\) 的一个子群。另一方面，回忆（第八章，§2，第 2 节，注 4）：Weyl 群 \(W(\mathfrak{g}_{\mathbf{C}}, t_{\mathbf{C}})\)（分裂约化代数（\(\mathfrak{g}_{\mathbf{C}}, t_{\mathbf{C}}\)）的 Weyl 群）作用在 \(t_{\mathbf{C}}\) 上，从而典则地等同于 \(\mathbf{GL}(t_{\mathbf{C}})\) 的一个子群。

命题 10. 映射 \(u \mapsto \mathrm{L}(u)_{(\mathbf{C})}\)（从 \(\operatorname{Aut}(\mathrm{T})\) 到 \(\mathbf{GL}(\mathfrak{t}_{\mathbf{C}})\)）诱导出从 \(\mathrm{W}\) 到分裂约化李代数 \((\mathfrak{g}_{\mathbf{C}}, \mathfrak{t}_{\mathbf{C}})\) 的 Weyl 群的一个同构。对所有 \(\alpha \in \mathrm{R}\)，\(\mathrm{W}_{\mathrm{Z}_{\alpha}}(\mathrm{T})\) 的阶为 2，且 \(\mathrm{W}_{\mathrm{Z}_{\alpha}}(\mathrm{T})\) 的非恒等元在上述同构下的像是反射 \(s_{H_{\delta(\alpha)}}\)。

所考虑的映射是单射。剩下要证明它的像等于 \(\mathrm{W}(\mathfrak{g}_{\mathbf{C}}, t_{\mathbf{C}})\)。

设 \(g \in \mathrm{N}_{\mathrm{G}}(\mathrm{T})\)。采用第八章 §5 第 2 节的记号，我们有 \(\operatorname{Ad} g \in \operatorname{Aut}(\mathfrak{g}_{\mathbf{C}}, \mathfrak{t}_{\mathbf{C}}) \cap \operatorname{Int}(\mathfrak{g}_{\mathbf{C}})\)，故 \(\operatorname{Ad} g \in \operatorname{Aut}_{0}(\mathfrak{g}_{\mathbf{C}}, \mathfrak{t}_{\mathbf{C}})\)（同前处，第 5 节，命题 11）。由同前处第 2 节命题 4，\(\mathfrak{t}_{\mathbf{C}}\) 上由 \(\operatorname{Ad} g\) 诱导的自同构属于 \(\mathrm{W}(\mathfrak{g}_{\mathbf{C}}, \mathfrak{t}_{\mathbf{C}})\)。于是，\(\mathrm{W}\) 在 \(\mathbf{GL}(\mathfrak{t}_{\mathbf{C}})\) 中的像含于 \(\mathrm{W}(\mathfrak{g}_{\mathbf{C}}, \mathfrak{t}_{\mathbf{C}})\)。

设 \(\alpha \in \mathrm{R}(\mathrm{G},\mathrm{T})\)，并设 \(\nu : \mathbf{SU}(2,\mathbf{C}) \to \mathrm{G}\) 是具有定理 1 的推论中所述性质的李群态射。\(\nu\) 在元素 \(\theta\)（\(\mathbf{SU}(2,\mathbf{C})\) 的元素）上的像具有以下性质（§3，第 6 节，公式 (17)）：

\[
a) (\operatorname{Int} \nu (\theta)) (t) = t \quad \text{ 若 } t \in \operatorname{Ker} \alpha ,
\]

\[
b) (\operatorname{Int} \nu (\theta)) (t) = t ^ {- 1} \quad \text{ 若 } t \in \mathrm{T} \cap \mathrm{S} _ {\alpha}.
\]

由此可知，\(\operatorname{Ad}\nu(\theta)\) 在 \(\operatorname{Ker}\delta(\alpha)\subset\mathfrak{t}_{\mathbf{C}}\) 上诱导出恒等映射，并诱导出映射 \(x\mapsto-x\)（在 \([\mathfrak{g}^{\alpha},\mathfrak{g}^{-\alpha}]\) 上），因而与反射 \(s_{H_{\delta(\alpha)}}\) 一致。于是，W 的像包含所有的 \(s_{H_{\delta(\alpha)}}\)，从而等于 \(\mathrm{W}(\mathfrak{g}_{\mathbf{C}}, t_{\mathbf{C}})\)。特别地，\(\mathrm{W}_{\mathrm{Z}_{\alpha}}(\mathrm{T})\) 的阶为 2，因而由恒等映射与 \(\operatorname{Int}\nu(\theta)\) 组成。命题证毕。

推论. 假设 G 是半单的。则 G 的每个元素都是 G 中两个元素的换位子。

设 \(c\) 是 Weyl 群 \(\mathrm{W}(\mathfrak{g}_{\mathbf{C}},\mathfrak{t}_{\mathbf{C}})\) 的一个 Coxeter 变换（第五章，§6，第 1 节），并设 \(n\) 是 \(\mathrm{N}_{\mathrm{G}}(\mathrm{T})\) 的一个元素，它在 \(\mathrm{W}\) 中的类经由命题中定义的同构等同于 \(c\)。用 \(f_{c}\) 表示态射 \(t\mapsto (n,t)\)（从 \(\mathrm{T}\) 到 \(\mathrm{T}\)）；对 \(x\in \mathfrak{t}_{\mathbf{C}}\)，我们有 \(\mathrm{L}(f_c)_{(\mathbf{C})}(x) = (\mathrm{Ad}n)(x) - x = c(x) - x\)。

由第五章 §6 第 2 节定理 1，自同态 \(c\)（\(\mathbf{t}_{\mathbf{C}}\) 上的）没有等于 1 的特征值。因此，\(\mathrm{L}(f_c)\) 是满射，从而 \(f_c\) 也是满射。由此可知，T 的每个元素都是 G 中两个元素的换位子，结合 §2 第 2 节定理 2 即得推论。

\subsection{基本群}

在下述命题中，\(f(\mathrm{G},\mathrm{T})\) 表示从 \(\varGamma(\mathrm{T})\) 到 \(\pi_{1}(\mathrm{G})\) 的同态，它是 \(\varGamma(\mathrm{T})\) 到 \(\pi_{1}(\mathrm{T})\) 的典则同构（第 2 节，注 3）与同态 \(\pi_{1}(\iota)\) 的复合，其中 \(\iota\) 是典则单射 \(T\to G\)。

命题 11. 同态 \(f(\mathrm{G},\mathrm{T}) : \Gamma(\mathrm{T}) \to \pi_1(\mathrm{G})\) 是满射。其核是 \(\mathrm{N}(\mathrm{G},\mathrm{T})\)（\(\Gamma(\mathrm{T})\) 的子群），由结点向量族 \((K_{\alpha})_{\alpha \in \mathrm{R}(\mathrm{G},\mathrm{T})}\) 生成。

同态 \(f(\mathrm{G},\mathrm{T})\) 的满射性由 \(\S2\) 第 4 节命题 3 得出。我们用 \(A(\mathrm{G},\mathrm{T})\) 表示断言：「\(f(\mathrm{G},\mathrm{T})\) 的核由诸 \(K_{\alpha}\) 生成」，它尚待证明；我们区分几种情形：

\begin{enumerate}
\def\labelenumi{\alph{enumi})}
\item
  G 是单连通的。设 \(\rho : \mathfrak{g}_{\mathbf{C}} \to \mathfrak{gl}(\mathrm{V})\) 是 \(\mathfrak{g}_{\mathbf{C}}\) 在有限维复向量空间 V 上的一个线性表示。限制到 \(\mathfrak{g}\)，就得到 \(\mathfrak{g}\) 在实向量空间 \(V_{(\mathbf{R})}\) 上的一个表示；由于 G 单连通，存在解析线性表示 \(\pi\)（G 在 \(V_{(\mathbf{R})}\) 上的）使得 \(\rho = L(\pi)\)。由第 3 节命题 7 可知，\(\delta(X(T))\)（X(T) 在 \(\mathfrak{t}_{\mathbf{C}}^{*}\) 中的像）包含 \(\rho\) 在 V 上的所有权。由于这对每个表示 \(\rho\)（\(\mathfrak{g}_{\mathbf{C}}\) 的表示）都成立，由第八章 §7 第 2 节定理 1 可知，\(\delta(X(T))\) 包含 \(\delta(R)\) 的权群，后者按定义是这样的集合：由 \(\lambda \in \mathfrak{t}_{\mathbf{C}}^{*}\) 中满足 \(\lambda(H_{\delta(\alpha)}) \in \mathbf{Z}\)（对所有 \(\alpha \in R\)）者组成，换言之，由满足 \(\lambda(K_{\alpha}) \in 2\pi i\mathbf{Z}\)（对所有 \(\alpha \in R\)）者组成。于是，群 X(T) 包含每个这样的元素 \(\lambda\)：它属于 X(T) \(\otimes\) Q，且 \(\langle\lambda, K_{\alpha} \rangle \in \mathbf{Z}\) 对所有 \(\alpha \in R\) 成立；由对偶性，这蕴含 \(\Gamma(T)\) 由诸 \(K_{\alpha}\) 生成，从而得到断言 \(A(G,T)\)。
\item
  G 是一个单连通群 \(G'\) 与一个环面 S 的直积。此时 T 是 \(T'\)（\(G'\) 的一个极大环面）与 S 的直积，\(\varGamma(\mathrm{T})\) 可等同于 \(\varGamma(\mathrm{T}')\times\varGamma(\mathrm{S})\)，\(\pi_{1}(\mathrm{G})\) 可等同于 \(\pi_{1}(\mathrm{G}')\times\pi_{1}(\mathrm{S})\)，而 \(f(\mathrm{G},\mathrm{T})\) 可等同于分量为 \(f(\mathrm{G}^{\prime},\mathrm{T}^{\prime})\) 与 \(f(\mathrm{S},\mathrm{S})\) 的同态。由于 \(f(\mathrm{S},\mathrm{S})\) 是双射，典则映射 \(\varGamma(\mathrm{T}^{\prime})\to\varGamma(\mathrm{T})\) 把 \(\operatorname{Ker}f(\mathrm{G}^{\prime},\mathrm{T}^{\prime})\) 双射地映到 \(\operatorname{Ker}f(\mathrm{G},\mathrm{T})\) 上。此外，诸 \(K_{\alpha}\) 属于 G 的换位子群 \(G^{\prime}\) 的李代数，从而属于 \(\varGamma(\mathrm{T}^{\prime})\) 的像，故显然 \(A(\mathrm{G}^{\prime},\mathrm{T}^{\prime})\) 蕴含 \(A(\mathrm{G},\mathrm{T})\)，结合 a) 即得断言 \(A(\mathrm{G},\mathrm{T})\)。
\item
  一般情形。存在核为有限的满射态射 \(p : G' \to G\)，其中 \(G'\) 是一个单连通群与一个环面的直积（\(\S1\)，第 4 节，命题 4）。若 \(T'\) 是 T 在 \(G'\) 中的逆像（这是 \(G'\) 的一个极大环面，见 \(\S2\)，第 3 节，命题 1），且 N 是 p 的核，则我们有正合序列 \(0 \to N \to G' \to G \to 0\) 与 \(0 \to N \to T' \to T \to 0\)，从而有一个各行均正合的交换图（第 2 节，注 1 及《一般拓扑》，第十一章，编写中）
\end{enumerate}

\[
\begin{array}{c c c c c c c c c} 0 & \longrightarrow & \Gamma (\mathrm{T} ^ {\prime}) & \longrightarrow & \Gamma (\mathrm{T}) & \longrightarrow & \mathrm{N} & \longrightarrow & 0 \\ & & \Big \downarrow f (\mathrm{G} ^ {\prime}, \mathrm{T} ^ {\prime}) & & \Big \downarrow f (\mathrm{G}, \mathrm{T}) & & \Big \downarrow \mathrm{Id} _ {\mathrm{N}} \\ 0 & \longrightarrow & \pi_ {1} (\mathrm{G} ^ {\prime}) & \longrightarrow & \pi_ {1} (\mathrm{G}) & \longrightarrow & \mathrm{N} & \longrightarrow & 0. \end{array}
\]

由蛇形图（《代数》，第 X 章，第 4 页，命题 2）立即可知，\(A(\mathrm{G}', \mathrm{T}')\) 蕴含 \(A(\mathrm{G}, \mathrm{T})\)，结合 \(b\)) 即得命题。

推论 1. G 是单连通的当且仅当族 \((K_{\alpha})_{\alpha \in \mathrm{R}(\mathrm{G},\mathrm{T})}\) 生成 \(\Gamma (\mathrm{T})\)。

推论 2. 设 H 是 G 的包含 T 的连通闭子群；则有正合序列

\[
0 \longrightarrow \mathrm{N} (\mathrm{H}, \mathrm{T}) \longrightarrow \mathrm{N} (\mathrm{G}, \mathrm{T}) \longrightarrow \pi_ {1} (\mathrm{H}) \longrightarrow \pi_ {1} (\mathrm{G}) \longrightarrow 0.
\]

这由《代数》第 X 章第 4 页命题 2（蛇形图）应用于下面的交换图得出

\[
\begin{array}{c c c c c c c c c} 0 & \longrightarrow & \mathrm{N(H,T)} & \longrightarrow & \varGamma (\mathrm{T}) & \longrightarrow & \pi_ {1} (\mathrm{H}) & \longrightarrow & 0 \\ & & \Big \downarrow & & \Big \downarrow & & \Big \downarrow & & \\ 0 & \longrightarrow & \mathrm{N(G,T)} & \longrightarrow & \varGamma (\mathrm{T}) & \longrightarrow & \pi_ {1} (\mathrm{G}) & \longrightarrow & 0. \end{array}
\]

注. 可以证明（参看 §5 习题 2）\(\pi_2(\mathrm{G / H})\) 为零。于是由前述序列的正合性，得到从 \(\pi_2(\mathrm{G / H})\) 到 \(\mathrm{N(G,T) / N(H,T)}\) 的一个同构。

推论 3. 同态 \(\pi_1(\mathrm{D}(\mathrm{G})) \to \pi_1(\mathrm{G})\)（对应于 \(\mathrm{D}(\mathrm{G})\) 到 \(\mathrm{G}\) 中的包含关系）诱导出从 \(\pi_1(\mathrm{D}(\mathrm{G}))\) 到 \(\pi_1(\mathrm{G})\) 的挠子群的同构。

实际上，\(\mathrm{T} \cap \mathrm{D}(\mathrm{G})\) 是 \(\mathrm{D}(\mathrm{G})\) 的一个极大环面（\(\S 2\)，第 3 节，命题 1 \(c\)）；由正合序列

\[
0 \longrightarrow \Gamma (\mathrm{T} \cap \mathrm{D} (\mathrm{G})) \longrightarrow \Gamma (\mathrm{T}) \longrightarrow \Gamma (\mathrm{T} / (\mathrm{T} \cap \mathrm{D} (\mathrm{G}))) \longrightarrow 0
\]

与命题 11，我们得到正合序列

\[
0 \longrightarrow \pi_ {1} (\mathrm{D} (\mathrm{G})) \longrightarrow \pi_ {1} (\mathrm{G}) \longrightarrow \Gamma (\mathrm{T} / (\mathrm{T} \cap \mathrm{D} (\mathrm{G}))) \longrightarrow 0,
\]

从而得到推论，因为 \(\pi_{1}(\mathrm{D}(\mathrm{G}))\) 有限而 \(\Gamma(\mathrm{T}/(\mathrm{T}\cap\mathrm{D}(\mathrm{G})))\) 自由。

\subsection{极大秩子群}

回忆（第六章，§1，第 7 节）：R = R(G, T) 的子集 P 称为闭的，如果 \((P + P) \cap R \subset P\)；称为对称的，如果 P = -P。

命题 12. 设 H 是 G 的包含 T 的连通闭子群全体所成的集合，按包含关系排序。映射 \(\mathrm{H} \mapsto \mathrm{R}(\mathrm{H}, \mathrm{T})\) 是从 H 到 \(\mathrm{R}(\mathrm{G}, \mathrm{T})\) 的对称闭子集全体所成的集合（按包含关系排序）上的一个递增双射。

若 \(H \in H\)，则 \(\mathrm{L}(\mathrm{H})_{(\mathbf{C})}\) 是 \(t_{C}\) 与诸 \(g^{\alpha}\)（\(\alpha \in \mathrm{R}(\mathrm{H}, \mathrm{T})\)）的直和；由于它是 \(g_{C}\) 中的一个约化子代数，故 R 的子集 \(\mathrm{R}(\mathrm{H}, \mathrm{T})\) 满足所述条件（第八章，§3，第 1 节，引理 2 与命题 2）。反之，若 P 是满足这些条件的 R 的子集，则 \(t_{C} \oplus \sum_{\alpha \in P} g^{\alpha}\) 是 \(g_{C}\) 的一个子代数（同前处），它在 R 上是有理的（第 3 节），因而形如 \(h_{(C)}\)，其中 h 是 g 的子代数。设 \(\mathrm{H}(\mathrm{P})\) 是由 h 定义的 G 的整子群；它是闭的（§2，第 4 节，注 1）。我们立即验证：映射 \(\mathrm{H} \mapsto \mathrm{R}(\mathrm{H}, \mathrm{T})\) 与 \(\mathrm{P} \mapsto \mathrm{H}(\mathrm{P})\) 都是递增的且互为逆映射。

推论 1. G 中包含 T 的闭子群只有有限多个。

设 H 是这样一个子群；则 \(H_{0} \in H\)，且 H 有限。此外，H 是 \(\mathrm{N}_{\mathrm{G}}(\mathrm{H}_{0})\) 的包含 \(H_{0}\) 的子群，而 \(\mathrm{N}_{\mathrm{G}}(\mathrm{H}_{0})/\mathrm{H}_{0}\) 有限（\(\S2\)，第 4 节，命题 4 与注 2）。

推论 2. 设 H 是 G 的包含 T 的连通闭子群，并设 \(W_{\mathrm{G}}^{\mathrm{H}}(\mathrm{T})\) 是 \(W_{\mathrm{G}}(\mathrm{T})\) 中 R 的子集 \(\mathrm{R}(\mathrm{H},\mathrm{T})\) 的稳定子群。群 \(N_{\mathrm{G}}(\mathrm{H})/\mathrm{H}\) 同构于商群 \(W_{\mathrm{G}}^{\mathrm{H}}(\mathrm{T})/W_{\mathrm{H}}(\mathrm{T})\)。

实际上，把 §2 第 5 节命题 7 应用于 \(\mathrm{N}_{\mathrm{G}}(\mathrm{H})\)，可知 \(\mathrm{N}_{\mathrm{G}}(\mathrm{H})/\mathrm{H}\) 同构于 \(\mathrm{W}_{\mathrm{N}(\mathrm{H})}(\mathrm{T})/\mathrm{W}_{\mathrm{H}}(\mathrm{T})\)，其中 \(\mathrm{W}_{\mathrm{N}(\mathrm{H})}(\mathrm{T})\) 是 \(\mathrm{W}_{\mathrm{G}}(\mathrm{T})\) 中其代表在 \(\mathrm{N}_{\mathrm{G}}(\mathrm{T})\) 中且正规化 H 的元素全体所成的集合。设 \(n \in \mathrm{N}_{\mathrm{G}}(\mathrm{T})\)，并设 w 是它在 \(\mathrm{W}_{\mathrm{G}}(\mathrm{T})\) 中的类。由第三章 §9 第 4 节命题 11，n 正规化 H 当且仅当 \((\operatorname{Ad} n)(\operatorname{L}(\operatorname{H})) = \operatorname{L}(\operatorname{H})\)；由第 3 节命题 5，这也意味着 R 的子集 \(\operatorname{R}(\operatorname{H}, \operatorname{T})\) 在 w 下稳定，由此得推论。

注 1. 群 \(W_{\mathrm{G}}^{\mathrm{H}}(\mathrm{T})\) 也是 \(W_{\mathrm{G}}(\mathrm{T})\) 中 T 的子群 \(\mathrm{C}(\mathrm{H})\) 的稳定子群：这由第 4 节命题 8 得出。

命题 13. 设 H 是 G 的一个极大秩连通闭子群，C 是它的中心。则 C 包含 G 的中心，且 H 是 C 的中心化子的单位元连通分支。

设 S 是 H 的一个极大环面。由于 G 的中心含于 S，故它含于 C。令 \(L = Z(C)_{0}\)；这是 G 的一个包含 H 的连通闭子群，故为极大秩，且其中心等于 C。用 \(R_{H}\) 与 \(R_{L}\) 分别表示 H 与 L 关于 S 的根系；则 \(R_{H} \subset R_{L} \subset R(G, S)\)。由于 \(C(H) = C(L)\)，命题 8（第 4 节）蕴含等式 \(Q(R_{H}) = Q(R_{L})\)；但 \(Q(R_{H}) \cap R_{L} = R_{H}\)（第六章，§1，第 7 节，命题 23），故 \(R_{H} = R_{L}\)，从而 H = L（命题 12）。

注 2. 称子群 \(\mathrm{C}\)（\(\mathrm{G}\) 中）为根子群，如果存在极大环面 \(\mathrm{S}\)（\(\mathrm{G}\) 中）与子集 \(\mathrm{P}\)（\(\mathrm{R}(\mathrm{G},\mathrm{S})\) 中），使得 \(\mathrm{C} = \bigcap_{\alpha \in \mathrm{P}}\mathrm{Ker}\alpha\)。由此

并结合命题 13 与第 5 节引理 2 可知，映射 \(H \mapsto C(H)\) 诱导出从极大秩连通闭子群全体所成的集合到 G 的根子群全体所成的集合的双射。其逆双射是映射 \(C \mapsto Z(C)_{0}\)。

推论. \(g \in \mathbf{G}\) 中满足 \(\mathrm{T} \cap g\mathrm{T}g^{-1} \neq \mathrm{C}(\mathrm{G})\) 者所成的集合，是 \(\mathbf{G}\) 中异于 \(\mathbf{G}\) 的有限多个闭解析子流形的并。

实际上，令 \(A_{g} = T \cap gTg^{-1}\)；我们有 \(T \subset Z(A_{g})\) 与 \(gTg^{-1} \subset Z(A_{g})\)。于是存在 \(x \in Z(A_{g})\) 使得 \(xTx^{-1} = gTg^{-1}\)（§2，第 2 节，定理 2），这蕴含 \(g \in Z(A_{g}).N_{G}(T)\)。用 A 表示 G 中包含 T 且异于 G 的闭子群所成的有限集合（推论 1），并令 \(X = \bigcup_{H \in A} H.N_{G}(T)\)；这是 G 的有限多个闭子流形的并，且异于 G。若 \(A_{g} \neq C(G)\)，则 \(Z(A_{g}) \in A\)，故 g 属于 X。反之，若 \(g \in H.N_{G}(T)\)（其中 \(H \in A\)），则 \(A_{g}\) 包含 \(C(H)\)，故 \(A_{g} \neq C(G)\)（命题 13）。

命题 14. 设 X 是 T 的一个子集，并设 \(R_{X}\) 是根 \(\alpha \in \mathrm{R}(\mathrm{G}, \mathrm{T})\) 中满足 \(\alpha(\mathrm{X}) = \{1\}\) 者所成的集合。群 \(\mathrm{Z}_{\mathrm{G}}(\mathrm{X}) / \mathrm{Z}_{\mathrm{G}}(\mathrm{X})_{0}\) 同构于 \(\mathrm{W}_{\mathrm{G}}(\mathrm{T})\) 中稳定 X 的子群关于由诸反射 \(s_{\alpha}\)（\(\alpha \in R_{X}\)）所生成的子群所得的商。

令 \(H = Z_{\mathrm{G}}(\mathrm{X})\)；由于 \(\mathrm{L}(\mathrm{H})_{(\mathbf{C})}\) 是 \(g_{C}\) 中被 \(\mathrm{Ad}(\mathrm{X})\) 固定的点所成的集合，它是 \(t_{C}\) 与诸 \(g^{\alpha}\)（满足 \(\alpha(\mathrm{X}) = \{1\}\) 者）的和。因此，\(\mathrm{R}(\mathrm{H}_{0}, \mathrm{T}) = \mathrm{R}_{\mathrm{X}}\)，故 \(\mathrm{W}_{\mathrm{H}_{0}}(\mathrm{T})\) 由诸反射 \(s_{\alpha}\)（\(\alpha \in R_{X}\)）生成。此时只需应用 §2 第 5 节命题 7。

我们将在下文（\(\S5\)，第 3 节，定理 1）看到：若 G 单连通且 X 仅由一点组成，则中心化子 \(Z(X)\) 是连通的。

\subsection{根图式}

定义 2. 根图式（或简称图式，如无混淆之虞）是一个三元组 \(D = (M, M_{0}, R)\)，其中：

\((RD_{0})\) M 是一个有限型的自由 Z-模，且子模 \(M_{0}\) 是 M 的直因子；

\((\mathrm{RD}_{\mathrm{I}})\) R 是 M 的一个有限子集；\(\mathrm{R} \cup \mathrm{M}_0\) 生成 Q-向量空间 \(\mathbf{Q} \otimes \mathrm{M}\)；

\((\mathrm{RD}_{\mathrm{II}})\) 对所有 \(\alpha \in \mathrm{R}\)，存在元素 \(\alpha^{\vee}\)（\(\mathrm{M}^{*} = \mathrm{Hom}_{\mathbf{Z}}(\mathrm{M}, \mathbf{Z})\) 中的），使得 \(\alpha^{\vee}(\mathrm{M}_{0}) = 0\)、\(\alpha^{\vee}(\alpha) = 2\)，且 M 的自同态 \(x \mapsto x - \alpha^{\vee}(x)\alpha\) 保持 R 不变。

由第六章 §1 第 1 节，对所有 \(\alpha \in \mathbb{R}\)，元素 \(\alpha^{\vee}\)（\(\mathrm{M}^*\) 中的）由 \(\alpha\) 唯一确定；我们用 \(s_{\alpha}\) 表示 M 的自同态 \(x \mapsto x - \alpha^{\vee}(x)\alpha\)。此外（同前处），\(\mathbf{Q}\)-向量空间 \(\mathbf{Q} \otimes \mathrm{M}\) 是 \(\mathbf{Q} \otimes \mathrm{M}_0\) 与由 R 生成的向量子空间 V(R) 的直和，且 R 是 V(R) 中的一个根系（同前处，定义 1）。

R 的元素称为根图式 D 的根，而元素 \(\alpha^{\vee}\)（\(M^{*}\) 中的）称为逆根。由 M 的诸自同构 \(s_{\alpha}\) 生成的群称为 D 的 Weyl 群，记作 \(\mathrm{W(D)}\)；\(\mathrm{W(D)}\) 的元素在 \(M_{0}\) 上诱导出恒等变换，而在 \(\mathrm{V(R)}\) 上诱导出根系 R 的 Weyl 群中的变换。

例. 1) 对每个有限型的自由 \(\mathbf{Z}\)-模 \(M\)，三元组 \((M, M, \varnothing)\) 是一个根图式。

\begin{enumerate}
\def\labelenumi{\arabic{enumi})}
\setcounter{enumi}{1}
\tightlist
\item
  若 \(D = (M, M_{0}, R)\) 是一个根图式，设 \(M_{0}^{*}\) 是 \(V(R)\) 在 \(M^{*}\) 中的正交补，并设 \(R^{\vee}\) 是 D 的逆根全体所成的集合。则 \(D^{\vee} = (M^{*}, M_{0}^{*}, R^{\vee})\) 是一个根图式，称为 D 的逆图式。对所有 \(\alpha \in R\)，对称 \(s_{\alpha^{\vee}}\)（\(M^{*}\) 的）是对称 \(s_{\alpha}\)（M 的）的反步自同构；映射 \(w \mapsto {}^{t} w^{-1}\) 是从 \(W(D)\) 到 \(W(D^{\vee})\) 的同构。此外，\(V(R^{\vee})\) 可自然地等同于 Q-向量空间 \(V(R)\) 的对偶，此时 \(R^{\vee}\) 等同于 R 的逆根系。
\end{enumerate}

若把 \(\mathbf{M}^*\) 的对偶等同于 M，则 \(\mathrm{D}^{\vee}\) 的逆图式等同于 D。

\begin{enumerate}
\def\labelenumi{\arabic{enumi})}
\setcounter{enumi}{2}
\item
  设 \((\mathfrak{g},\mathfrak{h})\) 是一个分裂约化 \(\mathbf{Q}\)-李代数，\(\mathrm{M} \subset \mathfrak{h}\) 是一个许可格（第八章，§12，第 6 节，定义 1）。设 \(\mathrm{M}_0\) 是 M 中与 \((\mathfrak{g},\mathfrak{h})\) 的根正交的子群，并设 \(\mathrm{R}^\vee\) 是诸 \(H_{\alpha}\)（\(\alpha \in \mathrm{R}(\mathfrak{g},\mathfrak{h})\)）所成的集合。则 \((\mathrm{M},\mathrm{M}_0,\mathrm{R}^\vee)\) 是一个根图式，而 \((\mathrm{M}^*,\mathrm{M}_0^*,\mathrm{R}(\mathfrak{g},\mathfrak{h}))\) 是其逆图式。
\item
  设 V 是 Q 上的一个向量空间，R 是 V 中的一个根系；用 P(R) 表示 R 的权群，用 Q(R) 表示 R 的根权格（第六章，§1，第 9 节）。则 (Q(R), 0, R) 与 (P(R), 0, R) 都是根图式。根图式 (M, M₀, S) 同构于形如 (Q(R), 0, R)（相应地 (P(R), 0, R)）的图式，当且仅当 M 由 S 生成（相应地 M* 由 S \(^{v}\) 生成）。
\end{enumerate}

对 P(R) 的每个包含 Q(R) 的子群 X，(X, 0, R) 是一个根图式；而且，在同构意义下，每个满足 M₀ = 0（换言之，S 生成 M 的一个有限指数子群）的图式 (M, M₀, S) 都由这种方式得到。

根图式 \((\mathrm{M},\mathrm{M}_0,\mathrm{R})\) 称为既约的，如果根系 \(\mathbf{R}\) 是既约的（换言之（第六章，§1，第 4 节），如果关系式 \(\alpha ,\beta \in \mathbb{R}\)、\(\lambda \in \mathbf{Z}\)、\(\beta = \lambda \alpha\) 蕴含 \(\lambda = 1\) 或 \(\lambda = -1\)）。例 1) 与 3) 中的图式都是既约的。

\subsection{紧李群与根系}

用前一节引入的术语，第 4 节与第 5 节的结果中的一个重要部分可以总结为下面的定理：

定理 2. a) (X(T), X(T/(T ∩ D(G)), R(G, T))) 是一个既约根图式；其 Weyl 群由诸 X(w)（w ∈ W）组成；群 X(C(G)) 同构于 X(T) 关于由 R(G, T) 生成的子群所得的商。

\begin{enumerate}
\def\labelenumi{\alph{enumi})}
\setcounter{enumi}{1}
\item
  \((\Gamma(\mathrm{T}),\Gamma(\mathrm{C}(\mathrm{G})_{0}),\mathrm{R}^{\vee}(\mathrm{G},\mathrm{T}))\) 是一个既约根图式；其 Weyl 群由诸 \(\Gamma(w)\)（\(w\in W\)）组成；群 \(\pi_{1}(\mathrm{G})\) 同构于 \(\Gamma(\mathrm{T})\) 关于由 \(\mathrm{R}^{\vee}(\mathrm{G},\mathrm{T})\) 生成的子群所得的商。
\item
  若把 Z-模 X(T) 与 \(\Gamma(T)\) 中的每一个都等同于另一个的对偶（第 2 节，命题 3），则前述两个根图式中的每一个都等同于另一个的逆图式。
\end{enumerate}

用 \(\mathrm{D}^{*}(\mathrm{G},\mathrm{T})\) 表示图式 \((\mathrm{X}(\mathrm{T}),\mathrm{X}(\mathrm{T}/(\mathrm{T}\cap\mathrm{D}(\mathrm{G})),\mathrm{R}(\mathrm{G},\mathrm{T})))\)，并用 \(\mathrm{D}_{*}(\mathrm{G},\mathrm{T})\) 表示图式 \((\varGamma(\mathrm{T}),\varGamma(\mathrm{C}(\mathrm{G})_{0}),\mathrm{R}^{\vee}(\mathrm{G},\mathrm{T}))\)；它们分别称为 G（关于 T）的反变图式与共变图式。

例. 1) 若 G 是秩 1 的半单群，则 \(\mathrm{D}^{*}(\mathrm{G},\mathrm{T})\) 与 \(\mathrm{D}_{*}(\mathrm{G},\mathrm{T})\) 必然同构于两个图式 \(\Delta_{2}=(\mathbf{Z},0,\{2,-2\})\)、\(\Delta_{1}=(\mathbf{Z},0,\{1,-1\})\) 之一。若 G 同构于 \(\mathbf{SU}(2,\mathbf{C})\)，则 \(\mathrm{D}_{*}(\mathrm{G},\mathrm{T})\) 同构于 \(\Delta_{1}\)（因为 G 单连通），故 \(\mathrm{D}^{*}(\mathrm{G},\mathrm{T})\) 同构于 \(\Delta_{2}\)。若 G 同构于 \(\mathbf{SO}(3,\mathbf{R})\)，则 \(\mathrm{D}^{*}(\mathrm{G},\mathrm{T})\) 同构于 \(\Delta_{1}\)（因为 \(\mathrm{C}(\mathrm{G})=\{1\})\)），故 \(\mathrm{D}_{*}(\mathrm{G},\mathrm{T})\) 同构于 \(\Delta_{2}\)。

\begin{enumerate}
\def\labelenumi{\arabic{enumi})}
\setcounter{enumi}{1}
\item
  若 G 与 \(G'\) 是两个连通紧李群，其极大环面分别为 T 与 \(T'\)，且若 \(\mathrm{D}^{*}(\mathrm{G},\mathrm{T})=(\mathrm{M},\mathrm{M}_{0},\mathrm{R})\)，\(\mathrm{D}^{*}(\mathrm{G}',\mathrm{T}')=(\mathrm{M}',\mathrm{M}_{0}',\mathrm{R}')\)，则 \(\mathrm{D}^{*}(\mathrm{G}\times\mathrm{G}',\mathrm{T}\times\mathrm{T}')\) 可等同于 \((\mathrm{M}\oplus\mathrm{M}',\mathrm{M}_{0}\oplus\mathrm{M}_{0}',\mathrm{R}\cup\mathrm{R}')\)。共变图式的情形类似。
\item
  设 N 是 T 的一个在 G 中为中心的闭子群，并设 \((M, M_{0}, R)\) 是 G 关于 T 的反变图式。则 G/N 关于 T/N 的反变图式可等同于 \((M', M_{0}', R)\)，其中 \(M'\) 是 M 中由诸 \(\lambda\)（满足 \(\lambda(N) = \{1\}\) 者）组成的子群，而 \(M_{0}' = M' \cap M_{0}\)。
\item
  类似地，设 N 是一个有限交换群，并设 \(\varphi : \pi_{1}(G) \to N\) 是一个满射同态。设 \(G'\) 是相伴于这个同态的 G 的覆盖；这是一个连通紧李群，N 是它的一个中心子群（《一般拓扑》，第十一章，编写中），且 G 可自然地等同于 \(G'/N\)。设 \(T'\) 是 \(G'\) 的作为 T 的逆像的极大环面。若 \((P, P_{0}, S)\) 是 G 关于 T 的共变图式，则 \(G'\) 关于 \(T'\) 的共变图式可等同于 \((P', P_{0}', S)\)，其中 \(P'\) 是复合同态 \(\varphi \circ f(G, T) : P \to N\) 的核（参看第 6 节，命题 11），而 \(P_{0}' = P_{0} \cap P'\)。
\end{enumerate}

注. 1) 设 \(\mathfrak{c}\) 是 \(\mathfrak{g}_{\mathbf{C}}\) 的中心；则 \(\mathfrak{c} = \mathrm{L}(\mathrm{C}(\mathrm{G}))_{(\mathbf{C})}\)。在 G 关于 T 的图式与分裂约化代数 \((\mathfrak{g}_{\mathbf{C}}, t_{\mathbf{C}})\) 的根系与逆根系之间，我们有下列关系：

\begin{enumerate}
\def\labelenumi{\alph{enumi})}
\tightlist
\item
  从 \(\mathbf{C} \otimes \Gamma(\mathrm{T})\) 到 \(\mathfrak{t}_{\mathbf{C}}\) 的典则同构诱导出从 \(\mathbf{C} \otimes \Gamma(\mathrm{C}(\mathrm{G})_0)\) 到 \(\mathfrak{c}\) 的双射，以及从 \(1 \otimes \mathrm{R}^{\vee}(\mathrm{G}, \mathrm{T})\) 到 \(2\pi i.\mathrm{R}^{\vee}(\mathfrak{g}_{\mathbf{C}}, \mathfrak{t}_{\mathbf{C}})\) 的双射。
\end{enumerate}

从 \(\mathbf{C}\otimes\varGamma(\mathrm{C}(\mathrm{G})_{0})\) 到 c 的双射，以及从 \(1\otimes\mathrm{R}^{\vee}(\mathrm{G},\mathrm{T})\) 到 \(2\pi i.\mathrm{R}^{\vee}(\mathfrak{g}_{\mathbf{C}},\mathfrak{t}_{\mathbf{C}})\) 的双射。b) 从 \(\mathbf{C}\otimes\mathrm{X}(\mathrm{T})\) 到 \(t_{C}^{*}\)（\(t_{C}\) 的对偶）的典则同构诱导出从 \(\mathbf{C}\otimes\mathrm{X}(\mathrm{T}/(\mathrm{T}\cap\mathrm{D}(\mathrm{G})))\) 到 \(t_{C}\cap\mathcal{D}(\mathfrak{g})_{C}\) 的正交补的双射，以及从 \(1\otimes\mathrm{R}(\mathrm{G},\mathrm{T})\) 到 \(\mathrm{R}(\mathfrak{g}_{\mathbf{C}},\mathfrak{t}_{\mathbf{C}})\) 的双射。

\begin{enumerate}
\def\labelenumi{\arabic{enumi})}
\setcounter{enumi}{1}
\tightlist
\item
  假设群 G 是半单的；用 R（相应地 \(R^{\vee}\)）表示根系 R(G, T)（相应地 \(R^{\vee}(G, T)\)），于是我们有包含关系
\end{enumerate}

\[
\mathrm{Q} (\mathrm{R}) \subset \mathrm{X} (\mathrm{T}) \subset \mathrm{P} (\mathrm{R}) \quad \mathrm{Q} (\mathrm{R} ^ {\vee}) \subset \Gamma (\mathrm{T}) \subset \mathrm{P} (\mathrm{R} ^ {\vee}).
\]

有限交换群 \(\mathrm{P}(\mathrm{R})/\mathrm{Q}(\mathrm{R})\) 与 \(\mathrm{P}(\mathrm{R}^{\vee})/\mathrm{Q}(\mathrm{R}^{\vee})\) 相互对偶（第六章，§1，第 9 节）；若用 \(M^{*}\) 表示有限交换群 M 的对偶群，则由前面的典则同构我们得到

\[
\begin{array}{l l} \Gamma (\mathrm{T}) / \mathrm{Q} (\mathrm{R} ^ {\vee}) \to \pi_ {1} (\mathrm{G}) & \mathrm{P} (\mathrm{R} ^ {\vee}) / \Gamma (\mathrm{T}) \to \mathrm{C} (\mathrm{G}) \\ \mathrm{P} (\mathrm{R}) / \mathrm{X} (\mathrm{T}) \to (\pi_ {1} (\mathrm{G})) ^ {\wedge} & \mathrm{X} (\mathrm{T}) / \mathrm{Q} (\mathrm{R}) \to (\mathrm{C} (\mathrm{G})) ^ {\wedge}. \end{array}
\]

特别地，\(\pi_{1}(\mathrm{G})\) 与 \(\mathrm{C}(\mathrm{G})\) 的阶的乘积等于 \(\mathrm{R}(\mathrm{G},\mathrm{T})\) 的连通指数 f（同前处）。

现在设 \(G'\) 是另一个连通紧李群，\(T'\) 是 \(G'\) 的一个极大环面。设 \(f: G \to G'\) 是一个李群同构，满足 \(f(T) = T'\)；并用 \(f_{T}\) 表示它所定义的从 T 到 \(T'\) 的同构。则 \(\mathrm{X}(f_{\mathrm{T}})\) 是从 \(\mathrm{D}^{*}(\mathrm{G}', \mathrm{T}')\) 到 \(\mathrm{D}^{*}(\mathrm{G}, \mathrm{T})\) 的同构，记作 \(\mathrm{D}^{*}(f)\)；而 \(\Gamma(f_{\mathrm{T}})\) 是从 \(\mathrm{D}_{*}(\mathrm{G}, \mathrm{T})\) 到 \(\mathrm{D}_{*}(\mathrm{G}', \mathrm{T}')\) 的同构，记作 \(\mathrm{D}_{*}(f)\)。若 \(t \in T\)，且令 \(g = f \circ \operatorname{Int} t = (\operatorname{Int} f(t)) \circ f\)，则 \(\mathrm{D}^{*}(g) = \mathrm{D}^{*}(f)\)，\(\mathrm{D}_{*}(g) = \mathrm{D}_{*}(f)\)。

命题 15. 设 \(\varphi\) 是从 \(\mathrm{D}^*(\mathrm{G}',\mathrm{T}')\) 到 \(\mathrm{D}^*(\mathrm{G},\mathrm{T})\)（相应地，从 \(\mathrm{D}_*(\mathrm{G},\mathrm{T})\) 到 \(\mathrm{D}_*(\mathrm{G}',\mathrm{T}')\)）的同构。则存在同构 \(f:\mathrm{G}\to \mathrm{G}'\)，使得 \(f(\mathrm{T}) = \mathrm{T}'\) 且 \(\varphi = \mathrm{D}^*(f)\)（相应地 \(\varphi = \mathrm{D}_*(f)\)）；若 \(f_{1}\) 与 \(f_{2}\) 是两个这样的同构，则存在元素 \(t\)（\(\mathrm{T}\) 中的），使得 \(f_{2} = f_{1}\circ \operatorname {Int}t\)。

第二个断言由第 4 节命题 9 立即得出；例如，我们对共变图式证明第一个断言。用 \(\mathfrak{g}'\)（相应地 \(\mathfrak{t}'\)）表示 \(G'\)（相应地 \(T'\)）的李代数，并用 \(\mathfrak{g}_{\mathbf{C}}'\)（相应地 \(\mathfrak{t}_{\mathbf{C}}'\)）表示其复化李代数。由第八章 §4 第 4 节定理 2 (i)，存在同构 \(\psi : \mathfrak{g}_{\mathbf{C}} \to \mathfrak{g}_{\mathbf{C}}'\)，它把 \(\mathfrak{t}_{\mathbf{C}}\) 映到 \(\mathfrak{t}_{\mathbf{C}}'\)，且在 \(\Gamma(T) \subset \mathfrak{t}_{\mathbf{C}}\) 上诱导出给定的同构 \(\varphi : \Gamma(T) \to \Gamma(T')\)。于是 \(\mathfrak{g}\) 与 \(\psi^{-1}(\mathfrak{g}')\) 是 \(\mathfrak{g}_{\mathbf{C}}\) 的两个紧形式，它们有一个公共的交 \(\mathfrak{t}\)（在 \(\mathfrak{t}_{\mathbf{C}}\) 中）；由 §3 第 2 节命题 3，存在内自同构 \(\theta\)（\(\mathfrak{g}_{\mathbf{C}}\) 的），它在 \(\mathfrak{t}_{\mathbf{C}}\) 上诱导出恒等映射，且满足 \(\theta(\mathfrak{g}) = \psi^{-1}(\mathfrak{g}')\)。把 \(\psi\) 换为 \(\psi \circ \theta\)，可以假设 \(\psi\) 把 \(\mathfrak{g}\) 映到 \(\mathfrak{g}'\)。进而，由第 2 节命题 4，存在唯一的态射 \(f_{\mathrm{T}} : T \to T'\) 使得 \(\Gamma(f_{\mathrm{T}}) = \varphi\)。于是 \(\psi\) 在 \(\mathfrak{t}\) 上的限制是 \(L(f_{\mathrm{T}})\)，且由 §2 第 6 节命题 8，存在唯一的态射 \(f : G \to G'\)，它诱导出 \(f_{\mathrm{T}}\)（在 \(T\) 上），以及 \(\psi\)（在 \(g_{C}\) 上）。把前面的论证应用于 \(\varphi^{-1}\) 与 \(\psi^{-1}\)，就得到 f 的一个逆态射，因而 f 是同构。于是 \(\mathrm{D}_{*}(f)=\Gamma(f_{\mathrm{T}})=\varphi\)，命题得证。

注意，若 T 与 \(T'\) 是 G 的两个极大环面，则图式 \(\mathrm{D}^{*}(\mathrm{G},\mathrm{T})\) 与 \(\mathrm{D}^{*}(\mathrm{G},\mathrm{T}')\) 同构（若 \(g \in G\) 满足 \(gTg^{-1} = T'\)，则 \(\operatorname{Int} g\) 是 G 到自身的同构，且把 T 映到 \(T'\)）。用 \(\mathrm{D}^{*}(\mathrm{G})\) 表示 \(\mathrm{D}^{*}(\mathrm{G},\mathrm{T})\) 的同构类（《集合论》，第 II 章，§6，第 2 节）；这是一个只依赖于 G 的根图式，称为 G 的反变图式。G 的共变图式 \(\mathrm{D}_{*}(\mathrm{G})\) 可类似地定义，于是我们得到：

推论. 两个连通紧李群 G 与 \(G'\) 同构，当且仅当图式 \(\mathrm{D}^{*}(\mathrm{G})\) 与 \(\mathrm{D}^{*}(\mathrm{G}')\)（相应地，\(\mathrm{D}_{*}(\mathrm{G})\) 与 \(\mathrm{D}_{*}(\mathrm{G}')\)）相等。

命题 16. 对每个既约根图式 D，存在一个连通紧李群 G，使得 \(\mathrm{D}^{*}(\mathrm{G})\)（相应地 \(\mathrm{D}_{*}(\mathrm{G})\)）同构于 D。

\begin{enumerate}
\def\labelenumi{\alph{enumi})}
\item
  必要时把 D 换成其逆图式，我们归结为构造 G 使 \(\mathrm{D}^{*}(\mathrm{G})\) 同构于 D。设 \(\mathrm{D} = (\mathrm{M}, \mathrm{M}_{0}, \mathrm{R})\)；则 \(Q \otimes M\) 是 \(Q \otimes M_{0}\) 与由 R 生成的向量子空间 \(\mathrm{V}(\mathrm{R})\) 的直和。此外，由于逆根在 M 上取整数值，M 在 \(\mathrm{V}(\mathrm{R})\) 上（沿 \(Q \otimes M_{0}\) 平行）的投影含于 R 的权群 \(\mathrm{P}(\mathrm{R})\)，故 M 是 \(\mathrm{M}_{0} \oplus \mathrm{P}(\mathrm{R})\) 的一个有限指数子群。用 \(D'\) 表示图式 \((\mathrm{M}_{0} \oplus \mathrm{P}(\mathrm{R}), \mathrm{M}_{0}, \mathrm{R})\)。
\item
  设 \(\mathfrak{a}\) 是一个复半单李代数，其典则根系同构于 \(R \subset \mathbf{C} \otimes V(R)\)（第八章，§4，第 3 节），并设 \(\mathfrak{g}_1\) 是 \(\mathfrak{a}\) 的一个紧实形式（\(\S 3\)，第 2 节，定理 1）。设 \(G_1\) 是一个单连通实李群，其李代数同构于 \(\mathfrak{g}_1\)；则 \(G_1\) 是紧的（\(\S 1\)，第 4 节，定理 1）。设 \(T_1\) 是 \(G_1\) 的一个极大环面。由定理 1，图式 \(D^*(G_1, T_1)\) 同构于 \((P(R), 0, R)\)。
\item
  设 \(T_{0}\) 是一个维数等于 \(M_{0}\) 的秩的环面；则 \(\mathrm{D}^{*}(\mathrm{T}_{0},\mathrm{T}_{0})\) 同构于 \((\mathrm{M}_{0},\mathrm{M}_{0},\varnothing)\)，故 \(\mathrm{D}^{*}(\mathrm{G}_{1}\times\mathrm{T}_{0},\mathrm{T}_{1}\times\mathrm{T}_{0})\) 同构于 \(D'\)（例 2）。
\item
  最后，设 N 是 \(T_{1} \times T_{0}\) 中与 M 正交的有限子群。令 \(\mathrm{G} = (\mathrm{G}_{1} \times \mathrm{T}_{0})/\mathrm{N}\)，\(\mathrm{T} = (\mathrm{T}_{1} \times \mathrm{T}_{0})/\mathrm{N}\)。则 G 是一个连通紧李群，T 是 G 的一个极大环面，且 D(G, T) 同构于 D（例 3）。
\end{enumerate}

概要. 连通紧李群按同构的分类于是归结为既约根图式的分类。连通紧半单李群对应于既约根图式 \((M, M_{0}, R)\)（满足 \(M_{0} = 0\) 者）；给出这样一个图式，等价于给出 Q 上向量空间 V 中的一个既约根系 R 以及 V 的一个满足 \(Q(R) \subset M \subset P(R)\) 的子群 M。

注 3. 设 \(\mathrm{T}'\) 是 G 的另一个极大环面，B（相应地 \(\mathrm{B}'\)）是根系 \(\mathrm{R}(\mathrm{G},\mathrm{T})\)（相应地 \(\mathrm{R}(\mathrm{G}',\mathrm{T}')\)）的一个基（第六章，§1，第 5 节，定义 2）。存在 \(g\in \mathrm{G}\) 使得 Int \(g\) 把 T 映到 \(\mathrm{T}'\) 且把 B 映到 \(\mathrm{B}'\)，且这些元素组成模 Int(T) 的唯一陪集（由于 T 与 \(\mathrm{T}'\) 共轭，可以假设 \(\mathrm{T} = \mathrm{T}'\)，于是只需应用第六章 §1 第 5 节注 4 与第 4 节命题 9）。由此可知，从 T 到 \(\mathrm{T}'\) 的同构（由 Int \(g\) 诱导）与 \(g\) 的选取无关；因而 \(\mathrm{D}_{*}(\mathrm{Int}g)\) 与 \(\mathrm{D}^{*}(\mathrm{Int}g)\) 亦然。把第八章 §5 第 3 节注 2 加以必要的改动后照搬过来，我们现在就可以定义 G 的典则极大环面、G 的典则共变与反变根图式，\ldots。

\subsection{连通紧李群的自同构}

用 \(\mathrm{Aut}(\mathrm{G})\) 表示 G 的自同构所成的李群（第三章，§10，第 2 节），并用 \(\mathrm{Aut}(\mathrm{G},\mathrm{T})\) 表示 \(\mathrm{Aut}(\mathrm{G})\) 中由元素 \(u\)（满足 \(u(\mathrm{T}) = \mathrm{T}\) 者）组成的闭子群。我们已看到（§1，第 4 节，命题 4 的推论 5），\(\mathrm{Aut}(\mathrm{G})\) 的单位元连通分支是内自同构子群 \(\mathrm{Int}(\mathrm{G})\)；用 \(\mathrm{Int}_{\mathrm{G}}(\mathrm{H})\) 表示 G 的子群 H 在 \(\mathrm{Int}(\mathrm{G})\) 中的像。

设 D 是 G 关于 T 的共变图式；用 \(\operatorname{Aut}(D)\) 表示它的自同构群，并用 \(\mathrm{W}(D)\) 表示它的 Weyl 群。映射 \(u \mapsto \mathrm{D}_{*}(u)\) 是从 \(\operatorname{Aut}(G, T)\) 到 \(\operatorname{Aut}(D)\) 的同态。第 9 节命题 15 立即给出：

命题 17. 同态 \(\operatorname{Aut}(\mathrm{G},\mathrm{T})\to \operatorname{Aut}(\mathrm{D})\) 是满射，其核为 \(\operatorname{Int}_{\mathrm{G}}(\mathrm{T})\)。

注意 \(\operatorname{Aut}(G,T) \cap \operatorname{Int}(G) = \operatorname{Int}_{G}(N_{G}(T))\)，且 \(\operatorname{Int}_{G}(N_{G}(T))\) 在 \(\operatorname{Aut}(D)\) 中的像是 \(W(D)\)（第 5 节，命题 10）。于是，命题 17 给出同构

\[
\operatorname{Aut} (\mathrm{G}, \mathrm{T}) / (\operatorname{Aut} (\mathrm{G}, \mathrm{T}) \cap \operatorname{Int} (\mathrm{G})) \rightarrow \operatorname{Aut} (\mathrm{D}) / \mathrm{W} (\mathrm{D}).
\]

此外，\(\operatorname{Aut}(\mathrm{G}) = \operatorname{Int}(\mathrm{G}).\operatorname{Aut}(\mathrm{G}, \mathrm{T})\)。实际上，若 u 属于 \(\operatorname{Aut}(\mathrm{G})\)，则 \(u(\mathrm{T})\) 是 T 的一个极大环面，故与 T 共轭，从而存在 G 的内自同构 v 使得 \(u(\mathrm{T}) = v(\mathrm{T})\)，换言之，\(v^{-1}u \in \operatorname{Aut}(\mathrm{G}, \mathrm{T})\)。由此可知，\(\operatorname{Aut}(\mathrm{G})/\operatorname{Int}(\mathrm{G})\) 可等同于 \(\operatorname{Aut}(\mathrm{G}, \mathrm{T})/(\operatorname{Aut}(\mathrm{G}, \mathrm{T}) \cap \operatorname{Int}(\mathrm{G}))\)，故由前述结果，我们有正合序列

\[
1 \to \operatorname{Int} (\mathrm{G}) \to \operatorname{Aut} (\mathrm{G}) \to \operatorname{Aut} (\mathrm{D}) / \mathrm{W} (\mathrm{D}) \to 1.\tag{16}
\]

因此：

命题 18. 群 \(\operatorname{Aut}(\mathrm{G}) / \operatorname{Int}(\mathrm{G})\) 同构于 \(\operatorname{Aut}(\mathrm{D}) / \operatorname{W}(\mathrm{D})\)。

特别地，假设 G 是半单的；这时群 \(\operatorname{Aut}(D)\) 可等同于 \(\mathrm{A}(\mathrm{R}(\mathrm{G},\mathrm{T}))\)（第六章，§1，第 1 节）中由满足 \(u(\mathrm{X}(\mathrm{T})) \subset \mathrm{X}(\mathrm{T})\) 的元素 u 组成的子群，而子群 \(\mathrm{W}(\mathrm{D})\) 可等同于 \(\mathrm{W}(\mathrm{R}(\mathrm{G},\mathrm{T}))\)。

推论. 若 G 是单连通的，或 C(G) 仅由单位元素组成，则群 \(\operatorname{Aut}(\mathrm{G})/\operatorname{Int}(\mathrm{G})\) 同构于 R(G,T) 的邓金图的自同构群。

这由前述结果及第六章 §4 第 2 节命题 1 的推论得出。

我们现在要证明扩张 (16) 容有截面。

对所有 \(\alpha\in\mathrm{R}(\mathrm{G},\mathrm{T})\)，用 \(\mathrm{V}(\alpha)\) 表示 g 的满足 \(\mathrm{V}(\alpha)_{(\mathbf{C})}=\mathfrak{g}^{\alpha}+\mathfrak{g}^{-\alpha}\) 的 2 维向量子空间；并用 K 表示相伴于 g 的 Killing 型的二次型。

定义 3. (G,T) 的一个标架是一个偶 \((\mathrm{B},(U_{\alpha})_{\alpha \in \mathrm{B}})\)，其中 B 是 \(\mathrm{R}(\mathrm{G},\mathrm{T})\) 的一个基（第六章，§1，第 5 节，定义 2），且对所有 \(\alpha \in \mathrm{B}\)，\(U_{\alpha}\) 是 \(\mathrm{V}(\alpha)\) 中满足 \(\mathrm{K}(U_{\alpha}) = -1\) 的元素。

G 的一个标架是指 G 的一个极大环面 T 连同 \((G, T)\) 的一个标架。

引理 3. 设 \(B_{0}\) 是 \(\mathrm{R}(\mathrm{G},\mathrm{T})\) 的一个基。群 \(\operatorname{Int}_{\mathrm{G}}(\mathrm{T})\) 单传递地作用于\((\mathrm{G},\mathrm{T})\) 的形如 \((\mathrm{B}_{0},(U_{\alpha})_{\alpha\in\mathrm{B}_{0}})\) 的标架所成的集合上。

对所有 \(\alpha\in B_{0}\)，用 \(\mathrm{K}(\alpha)\) 表示二次型 K 在 \(\mathrm{V}(\alpha)\) 上的限制；T 在 \(\mathrm{V}(\alpha)\) 上的作用定义一个态射 \(\iota_{\alpha}:T\to\mathbf{SO}(\mathrm{K}(\alpha))\)。我们在第 4 节已看到，\(\mathbf{SO}(\mathrm{K}(\alpha))\) 可等同于 U，使得 \(\iota_{\alpha}\) 恰好等同于根 \(\alpha\)。由于 \(B_{0}\) 是 R 的基，它是由根生成的 Z-模 \(\mathrm{Q}(\mathrm{R})\) 的基，因而是 \(\mathrm{X}(\mathrm{T}/\mathrm{C}(\mathrm{G}))\)（\(\mathrm{X}(\mathrm{T})\) 的一个子模）的基。由此可知，诸态射 \(\iota_{\alpha}\) 的乘积诱导出从 \(\mathrm{T}/\mathrm{C}(\mathrm{G})\) 到诸群 \(\mathbf{SO}(\mathrm{K}(\alpha))\) 的乘积的同构。而后一群单传递地作用于第一分量为 \(B_{0}\) 的 (G,T) 的标架所成的集合上。

命题 19. 群 \(\operatorname{Int}(\mathbf{G})\) 单传递地作用于 \(\mathbf{G}\) 的标架所成的集合上。

设 \(e = (\mathrm{T}, \mathrm{B}, (U_{\alpha}))\) 与 \(e' = (\mathrm{T}', \mathrm{B}', (U_{\alpha}')\) 是 \(G\) 的两个标架。存在元素 \(g\) 属于 \(G\)，使得 \((\operatorname{Int} g)(\mathrm{T}) = \mathrm{T}'\)，且这些元素组成模 \(N_G(T)\) 的单个陪集。于是可以假设 \(T = T'\)，而我们必须证明存在 \(\operatorname{Int}_G(N_G(T))\) 的唯一元素把 \(e\) 变到 \(e'\)。由第六章 §1 第 5 节注 4，存在唯一的 \(w\) 属于 \(W(R)\)，使得 \(w(B) = B'\)。由于 \(W(R)\) 可等同于 \(N_G(T)/T\)，故存在 \(n \in N_G(T)\) 使得 \(w = \operatorname{Int} n\)，且 \(n\) 模 \(T\) 唯一确定。于是可以假设 \(B = B'\)，而我们必须证明存在 \(\operatorname{Int}_{\mathrm{G}}(\mathrm{T})\) 的唯一元素把 e 变到 \(e'\)，这正是引理 3。

推论. 设 e 是 (G,T) 的一个标架，并设 E 是使 e 稳定的 G 的自同构群。则 \(\operatorname{Aut}(G)\) 是 E 与 \(\operatorname{Int}(G)\) 的半直积，而 \(\operatorname{Aut}(G,T)\) 是 E 与 \(\operatorname{Int}(G)\cap\operatorname{Aut}(G,T)=\operatorname{Int}_{G}(\mathrm{N}_{\mathrm{G}}(\mathrm{T}))\) 的半直积。

实际上，\(\operatorname{Aut}(G)\) 的每个元素把 e 变为 G 的一个标架。由命题 19，\(\operatorname{Aut}(G)\) 模 \(\operatorname{Int}(G)\) 的每个陪集与 E 恰交于一点，由此得第一个断言。第二个断言同理可证。

注. 设 G 与 \(G'\) 是两个连通紧李群，并设 \(e = (\mathrm{T}, \mathrm{B}, (U_{\alpha}))\) 与 \(e' = (\mathrm{T}', \mathrm{B}', (U_{\alpha}')\) 分别是 G 与 \(G'\) 的标架。设 X 是从 G 到 \(G'\) 且把 e 变到 \(e'\) 的同构全体所成的集合。映射 \(f \mapsto \mathrm{D}^{*}(f)\)（相应地 \(\mathrm{D}_{*}(f)\)）是从 X 到这样的同构全体所成集合的双射：这些同构是从 \(\mathrm{D}^{*}(\mathrm{G}', \mathrm{T}')\) 到 \(\mathrm{D}^{*}(\mathrm{G}, \mathrm{T})\)（相应地，从 \(\mathrm{D}_{*}(\mathrm{G}, \mathrm{T})\) 到 \(\mathrm{D}_{*}(\mathrm{G}', \mathrm{T}')\)）的同构，且把 \(B'\) 变到 B（相应地，把 B 变到 \(B'\)）。实际上，这由命题 15 与引理 3 立即得出。

\section{共轭类}

我们沿用 §4 的记号。

\subsection{正则元}

由 §2 第 2 节定理 2 的推论 4，正则元 \(g\)（\(G\) 的）可由下列性质之一刻画：

\begin{enumerate}
\def\labelenumi{\alph{enumi})}
\item
  g 中被 Ad g 固定的子代数是一个嘉当子代数。
\item
  \(Z(g)_{0}\) 是 G 的一个极大环面。
\end{enumerate}

\(\mathrm{G}\) 的正则元所成的集合在 \(\mathrm{G}\) 中是开且稠密的。

在本段余下部分，我们用 \(G_{r}\)（相应地 \(T_{r}\)）表示 G（相应地 T）中在 G 中正则的点所成的集合。G 的元素 g 属于 \(T_{r}\) 当且仅当 \(Z(g)_{0}\) 等于 T；\(G_{r}\) 的每个元素都共轭于 \(T_{r}\) 的某个元素（§2，第 2 节，定理 2）。

T 的元素 t 属于 \(T_{r}\) 当且仅当 \(t^{\alpha} \neq 1\) 对每个根 \(\alpha \in \mathrm{R}(\mathrm{G}, \mathrm{T})\) 成立；因此，\(T - T_{r}\) 是诸子环面 Ker \(\alpha\)（\(\alpha\) 取遍 \(\mathrm{R}(\mathrm{G}, \mathrm{T})\)）的并。

命题 1. 令 \(n = \dim G\)。存在一个紧实解析流形 \(V\)（维数为 \(n - 3\)），以及一个解析映射 \(\varphi : V \to G\)，其像为 \(G - G_r\)。

设 \(\alpha\in\mathrm{R}(\mathrm{G},\mathrm{T})\)；令 \(\mathrm{V}_{\alpha}=(\mathrm{G}/\mathrm{Z}(\mathrm{Ker}\alpha))\times(\mathrm{Ker}\alpha)\)，并用 \(\varphi_{\alpha}\) 表示从 \(V_{\alpha}\) 到 G 的态射，使得对所有 \(g\in G\) 与所有 \(t\in\operatorname{Ker}\alpha\)，有 \(\varphi_{\alpha}(\bar{g},t) = gtg^{-1}\)（用 \(\bar{g}\) 表示 \(g\) 模 \(\mathrm{Z}(\mathrm{Ker}\alpha)\) 的陪集）。则 \(\mathrm{V}_{\alpha}\) 是一个维数为

\[
\begin{array}{c} \dim V _ {\alpha} = \dim G - \dim Z (\operatorname{Ker} \alpha) + \dim \operatorname{Ker} \alpha \\ = n - (\dim T + 2) + (\dim T - 1) = n - 3 \end{array}
\]

的紧实解析流形（§4，第 5 节，定理 1）；\(\varphi_{\alpha}\) 是实解析流形的态射，且 \(\varphi_{\alpha}\) 的像由 G 中共轭于 \(\mathrm{Ker} \alpha\) 中元素的元素组成。此时只需取 V 为诸流形 \(V_{\alpha}\) 之和，并取 \(\varphi\) 为诱导 \(\varphi_{\alpha}\)（在每个 \(V_{\alpha}\) 上）的态射。

注. 称 G 的元素 g 为强正则的，如果 \(Z(g)\) 是 G 的一个极大环面。若 \(g \in T\)，则 g 强正则当且仅当 \(w(g) \neq g\) 对 \(W_{G}(T)\) 的每个非恒等元素 w 成立（§4，第 7 节，命题 14）。于是，强正则元所成的集合是 G 的一个稠密开子集（§2，第 5 节，命题 5 的推论 2）。

\subsection{房与单房}

用 \(t_{r}\) 表示 t 中使 exp x 为正则的元素 \(x \in t\) 所成的集合，换言之，属于 \(T_{r}\) 的元素所成的集合。t 的元素 x 属于 \(t - t_{r}\) 当且仅当存在根 \(\alpha \in \mathrm{R}(\mathrm{G}, \mathrm{T})\) 使得 \(\delta(\alpha)(x) \in 2\pi i\mathbf{Z}\)。对每个根 \(\alpha \in \mathrm{R}(\mathrm{G}, \mathrm{T})\) 与每个整数 n，用 \(H_{\alpha,n}\) 表示诸 \(x \in t\) 中满足 \(\delta(\alpha)(x) = 2\pi in\) 者所成的集合。诸 \(H_{\alpha,n}\) 称为 t 的奇异超平面，而 \(t - t_{r}\) 是诸奇异超平面之并。t 的单房是 \(t_{r}\) 的连通分支，而房是 t 中过原点的那些奇异超平面（即 \(H_{\alpha,0} = \operatorname{Ker}\delta(\alpha), \alpha \in \mathrm{R}(\mathrm{G}, \mathrm{T})\)）之并的补集的连通分支。

我们有 \(\varGamma(\mathrm{T})\subset\mathfrak{t}-\mathfrak{t}_{r}\)；用 N(G,T) 表示由结点向量生成的 \(\varGamma(\mathrm{T})\) 的子群（§4，第 5 节）；由 §4 第 6 节命题 11，商 \(\varGamma(\mathrm{T})/\mathrm{N}(\mathrm{G},\mathrm{T})\) 可等同于 G 的基本群。

最后，用 W 表示 G 关于 T 的 Weyl 群，把它视为 T 与 t 的自同构群，并用 \(W_{a}\)（相应地 \(W_{a}^{\prime}\)）表示由 W 与诸平移 \(t_{\gamma}: x \mapsto x + \gamma\)（\(\gamma \in \mathrm{N}(\mathrm{G}, \mathrm{T})\)，相应地 \(\gamma \in \Gamma(\mathrm{T})\)）所生成的仿射空间 t 的自同构群。

设 \(w \in W\)，\(\gamma \in \Gamma(T)\)，\(\alpha \in R(G, T)\) 且 \(n \in Z\)。我们有：

\[
w (\mathrm{H} _ {\alpha , n}) = \mathrm{H} _ {w \alpha , n}, t _ {\gamma} (\mathrm{H} _ {\alpha , n}) = \mathrm{H} _ {\alpha , n + \langle \gamma , \alpha \rangle}.
\]

由此可知，对所有房 C 与所有 \(w \in W\)，\(w(\mathrm{C})\) 是一个房，且对所有单房 A 与所有 \(w \in W_{a}^{\prime}\)，\(w(\mathrm{A})\) 是一个单房。注意，当 \(\mathrm{X}(\mathrm{T}) \otimes \mathbf{R}\) 等同于 \(t^{*}\)（借助同构 \((2\pi i)^{-1}\delta\)）时，t 的诸单房与群 \(W_{a}\) 就是相伴于根系 \(\mathrm{R}(\mathrm{G}, \mathrm{T})\) 的诸单房与仿射 Weyl 群（第六章，§2，第 1 节）。

命题 2. a) 群 \(W_{a}\)（相应地 \(W_{a}^{\prime}\)）是 W 与 N(G, T)（相应地 \(\Gamma(T)\)）的半直积；子群 \(W_{a}\)（\(W_{a}^{\prime}\) 的）是正规的。

\begin{enumerate}
\def\labelenumi{\alph{enumi})}
\setcounter{enumi}{1}
\item
  群 W（相应地 \(\mathrm{W}_a\)）单传递地作用于诸房（相应地诸单房）所成的集合上。
\item
  设 C 是一个房，A 是一个单房。则 \(\overline{C}\)（相应地 \(\overline{A}\)，相应地 A）是 W 在 t 上的作用（相应地 \(W_{a}\) 在 t 上的作用，相应地 \(W_{a}\) 在 \(t-t_{r}\) 上的作用）的基本域。若 \(x \in t_{r}\) 与 \(w \in W_{a}\) 满足 \(w(x) = x\)，则 w = Id。
\item
  对每个房 C，存在唯一的单房 A 使得 \(A \subset C\) 且 \(0 \in \overline{A}\)。对每个单房 A，存在唯一的 \(\gamma \in \mathrm{N}(\mathrm{G}, \mathrm{T})\) 使得 \(\gamma \in \overline{A}\)。
\end{enumerate}

若 \(w \in W\) 且 \(\gamma \in \Gamma(T)\)，则 \(wt_{\gamma}w^{-1} = t_{w(\gamma)}\)，且 \(wt_{\gamma}w^{-1}t_{\gamma}^{-1} = t_{w(\gamma)-\gamma}\)，其中 \(w(\gamma) - \gamma \in \mathrm{N}(\mathrm{G}, \mathrm{T})\)；由此立得 a)。命题的其余部分由第六章 §1 第 5 节与 §2 第 1、2 节得出。

推论 1. 设 A 是 t 的一个单房，\(\overline{A}\) 是它的闭包，\(H_{A}\) 是 A 在 \(W_{a}^{\prime}\) 中的稳定子群。

\begin{enumerate}
\def\labelenumi{\alph{enumi})}
\item
  群 \(W_{a}^{\prime}\) 是 \(H_{A}\) 与 \(W_{a}\) 的半直积。
\item
  指数映射 \(\overline{\mathrm{A}}\to \mathrm{T}\) 与典则单射 \(\mathrm{T}\rightarrow \mathrm{G}\) 通过过渡到商与子集，诱导出同胚
\end{enumerate}

\[
\begin{array}{l} \overline {{\mathrm{A}}} / \mathrm {H_ {A}} \to \mathrm{T} / \mathrm{W} \to \mathrm{G} / \mathrm{Int} (\mathrm{G}) \\ \mathrm{A} / \mathrm {H_ {A}} \to \mathrm {T_ {r}} / \mathrm{W} \to \mathrm {G_ {r}} / \mathrm{Int} (\mathrm{G}). \end{array}
\]

设 \(w' \in W_a'\)；则 \(w'(A)\) 是 \(t\) 的一个单房，且存在（命题 2 b)）唯一的 \(w\) 属于 \(W_a\)，使得 \(w(A) = w'(A)\)，即 \(w^{-1}w' \in H_A\)。由于 \(W_a\) 在 \(W_a'\) 中正规，这就证明了 \(a\))。

\(\overline{A}\) 到 t 中的典则单射诱导出一个连续双射 \(\theta: \overline{A} \to t/W_{a}\)（命题 2 c)），而它是同胚，因为 \(\overline{A}\) 是紧的。由于 \(W_{a}\) 在 \(W_{a}^{\prime}\) 中正规，群 \(H_{A}\) 典则地作用于 \(t/W_{a}\) 上（《代数》，第一章，§5，第 4 节），且 \(t/W_{a}^{\prime}\) 可等同于商 \((t/W_{a})/H_{A}\)；映射 \(\theta\) 与 \(H_{A}\) 的作用相容，故通过过渡到商诱导出同胚 \(\overline{A}/H_{A} \to t/W_{a}^{\prime}\)。进而，\(\exp_{\Gamma}\) 诱导出从 \(t/\Gamma(T)\) 到 T 的一个同胚，从而也诱导出从 \(t/W_{a}^{\prime}\) 到 T/W 的一个同胚。断言 b) 由这一点以及 §2 第 4 节命题 5 的推论 1 得出。

注. 1) 群 \(H_{A}\) 可自然地等同于 \(\Gamma(\mathrm{T})/\mathrm{N}(\mathrm{G},\mathrm{T})\)，从而也可等同于 \(\pi_{1}(\mathrm{G})\)。于是，当 G 单连通时，它就化为单位元素。

\begin{enumerate}
\def\labelenumi{\arabic{enumi})}
\setcounter{enumi}{1}
\item
  设 \(x \in \mathrm{A}\)；则 \(\exp x \in \mathrm{T}_r\)，故 \(\mathrm{Z}(\exp x)_0 = \mathrm{T}\)。进而，\(\exp x\) 是强正则的（第 1 节，注）当且仅当 \(w(x) \neq x\) 对所有异于恒等映射的 \(w \in \mathrm{W}_a'\) 成立。由推论 1，这也意味着 \(h(x) \neq x\) 对所有异于恒等映射的 \(h \in \mathrm{H}_{\mathrm{A}}\) 成立。特别地，若 \(\mathrm{G}\) 单连通，则 \(\mathrm{Z}_{\mathrm{G}}(t) = \mathrm{T}\) 对所有 \(t \in \mathrm{T}_r\) 成立，且 \(\mathrm{G}\) 的每个正则元都是强正则的。
\item
  \(W_{a}\) 的特殊点（第六章，§2，第 2 节）是 t 中满足 \(\delta(\alpha)(x) \in 2\pi i\mathbf{Z}\)（对所有 \(\alpha \in \mathrm{R}(\mathrm{G}, \mathrm{T})\)）的元素 x（同前处，命题 3），也就是满足 \(\exp x \in \mathrm{C}(\mathrm{G})\) 的元素 x（§4，第 4 节，命题 8）。对这样的元素 x，我们有 \(wx - x \in \mathrm{N}(\mathrm{G}, \mathrm{T})\)（对所有 \(w \in W\)）（第六章，§1，第 9 节，命题 27），故 x 在 \(W_{a}\) 与 \(W_{a}^{\prime}\) 中的稳定子群一致。设 S 是 \(\overline{A}\) 的特殊点所成的集合；由前述结果与推论 1 可知，群 \(H_{A}\) 自由地作用于 S 上，且指数映射诱导出从 \(S/H_{A}\) 到 \(C(G)\) 的一个双射。
\end{enumerate}

推论 2. 设 C 是 t 的一个房，\(\overline{C}\) 是它的闭包。典则单射 \(\overline{C} \rightarrow t \rightarrow g\) 通过过渡到商诱导出同胚

\[
\overline {{\mathrm{C}}} \rightarrow \mathfrak {t} / \mathrm{W} \rightarrow \mathfrak {g} / \operatorname{Ad} (\mathrm{G}).
\]

典则映射 \(\overline{C} \rightarrow t\) 与 \(t \rightarrow t/W\) 是真的（《一般拓扑》，第三章，§4，第 1 节，命题 2 c)）。映射 \(\overline{C} \rightarrow t/W\) 连续、真且双射（命题 2 c)）；因而是同胚，于是结合 §2 第 5 节命题 6 的推论即得推论。

注 4. 用 \(\mathfrak{g}_{\mathrm{reg}}\) 表示 \(\mathfrak{g}\) 的正则元所成的集合（第七章，§2，第 2 节，定义 2），并令 \(\mathfrak{t}_{\mathrm{reg}} = \mathfrak{t} \cap \mathfrak{g}_{\mathrm{reg}}\)。对 \(x \in \mathfrak{t}\)，我们有

\[
\det (\mathrm{X} - \operatorname{ad} _ {\mathfrak {g}} x) = \mathrm{X} ^ {\dim t} \prod_ {\alpha \in \mathrm{R(G,T)}} (\mathrm{X} - \delta (\alpha) (x)),
\]

因而 \(t_{reg}\) 是 t 中满足 \(\delta(\alpha)(x) \neq 0\)（对所有 \(\alpha \in \mathrm{R}(\mathrm{G}, \mathrm{T})\)）的元素 x 所成的集合，也就是 t 的诸房之并（于是 \(t_r \subset t_{reg}\)）。由此，\(\overline{C} \cap t_{reg} = C\)，故我们有同胚

\[
\mathrm{C} \rightarrow \mathfrak {t} _ {\text { reg }} / \mathrm{W} \rightarrow \mathfrak {g} _ {\text { reg }} / \mathrm{Ad} (\mathrm{G}).
\]

推论 3. 假设 G 是单连通的；设 g 是 G 的一个正则元。存在 G 的一个极大环面 S 与 L(S) 的一个单房 A，二者均唯一确定，使得 \(g \in \exp(A)\) 且 \(0 \in \overline{A}\)。

可以假设 \(g\) 属于 \(\mathrm{T}_r\)（§2，第 2 节，定理 2）。设 \(x\) 是 \(\mathfrak{t}_r\) 中满足 \(\exp x = g\) 的一个元素，并设 \(\mathrm{A}'\) 是 \(\mathfrak{t}\) 中包含 \(x\) 的单房。诸单房 \(\mathrm{A}\)（\(\mathfrak{t}\) 中的）若满足 \(g \in \exp(\mathrm{A})\)，便是诸单房 \(\mathrm{A}' - \gamma\)（\(\gamma \in \Gamma(\mathrm{T})\)）；于是，断言由命题 2 d) 得出。

\subsection{自同构与正则元}

引理 1. 设 \(u\) 是 \(G\) 的一个自同构，\(H\) 是它的不动点所成的集合。a) \(H\) 是 \(G\) 的一个闭子群。

\begin{enumerate}
\def\labelenumi{\alph{enumi})}
\setcounter{enumi}{1}
\tightlist
\item
  若 \(\mathrm{H}_0\) 在 G 中是中心的，则 G 是交换的（于是 \(\mathrm{G} = \mathrm{T}\)）。
\end{enumerate}

断言 \(a\)) 显然。为证 \(b\))，我们可以把 \(\mathrm{G}\) 换成 \(\mathrm{D}(\mathrm{G})\)（\(\S 1\)，第 4 节，命题 4 的推论 1），从而可假设 \(\mathrm{G}\) 是半单的。这时，若 \(\mathrm{H}_0\) 在 \(\mathrm{G}\) 中是中心的，则我们有 \(\mathrm{L}(\mathrm{H}) = \{0\}\)，故自同态 \(\mathrm{L}(u) - \mathrm{Id}\)（\(\mathfrak{g}\) 上的）是双射的。设 \(f\) 是流形 \(\mathrm{G}\) 的由 \(f(g) = u(g)^{-1}g\)（对 \(g \in \mathrm{G}\)）定义的自同态；它是平展的，因为若 \(g \in \mathrm{G}\) 且 \(x \in \mathfrak{g}\)，我们有 \(\mathrm{T}(f)(xg) = u(g)^{-1}(x - \mathrm{L}(u)(x))g\)，故 \(f\) 在 \(g\) 处的切映射是双射。由此可知，\(f\) 的像是开且紧的，由于 \(\mathrm{G}\) 连通，故它与 \(\mathrm{G}\) 重合。现在设 \(\mathrm{E}\) 是 G 的一个标架（§4，第 10 节，定义 3），并设 \(u(\mathrm{E})\) 是它在 \(u\) 下的像。由同前处命题 19，存在 G 的元素 \(h\) 使得 \((\operatorname{Int} h)(\mathrm{E}) = u(\mathrm{E})\)。设 \(g \in \mathrm{G}\) 满足 \(h = f(g) = u(g)^{-1}g\)；则

\[
u \circ \operatorname{Int} g = (\operatorname{Int} u (g)) \circ u = \operatorname{Int} g \circ (\operatorname{Int} h) ^ {- 1} \circ u,
\]

于是标架 \((\operatorname{Int} g)(\operatorname{E})\) 在 u 下稳定。若 \((\operatorname{Int} g)(\operatorname{E}) = (\operatorname{T}_{1}, \operatorname{B}, (U_{\alpha})_{\alpha \in \operatorname{B}})\)，则 \(\sum U_{\alpha} \in \operatorname{L}(\operatorname{H})\)；由于 \(\operatorname{L}(\operatorname{H}) = \{0\}\)，这蕴含 \(B = \varnothing\)，故 \(G = T_{1}\)，从而 G 是交换的。

引理 2. 设 x 是 T 的一个元素，S 是 T 的一个子环面。若 \(Z(x) \cap Z(S)\) 的单位元连通分支就是 T，则存在 S 的元素 s 使得 xs 是正则的。

对所有 \(\alpha\in\mathrm{R}(\mathrm{G},\mathrm{T})\)，用 \(S_{\alpha}\) 表示 S 中由满足 \((xs)^{\alpha}=1\) 的元素 s 所组成的子流形。若不存在 S 的元素 s 使 xs 正则，则 S 是诸子流形 \(S_{\alpha}\) 之并，故等于其中之一。于是存在 \(\alpha\)（\(\mathrm{R}(\mathrm{G},\mathrm{T})\) 中）使得 \((xs)^{\alpha}=1\) 对所有 \(s\in S\) 成立；但这蕴含 \(x^{\alpha}=1\) 且 \(\alpha|S=1\)，故 \(\mathrm{Z}(x)\cap\mathrm{Z}(\mathrm{S})\supset\mathrm{Z}(\mathrm{Ker}\alpha)\)，引理得证。

引理 3. 假设 G 是单连通的。设 C 是 t 的一个房，u 是 G 的一个使 T 与 C 都稳定的自同构。则 T 中被 u 固定的点所成的集合是连通的。

由于 G 单连通，\(\varGamma(\mathrm{T})\) 由诸结点向量 \(K_{\alpha}\)（\(\alpha \in \mathrm{R}(\mathrm{G}, \mathrm{T})\)）生成，故以诸 \(K_{\alpha}\)（其中的 \(\alpha\) 属于由 C 定义的基 \(\mathrm{B}(\mathrm{C})\)）所成的族为一个基（第六章，§1，第 10 节）。于是只需证明：若 \(\varphi\) 是环面 T 的一个使 \(\varGamma(\mathrm{T})\) 的某个基保持稳定的自同构，则 \(\varphi\) 的不动点所成的集合是连通的。把这个基分解成由 \(\varphi\) 生成的群的诸轨道的不交并，我们就归结为 \(T = U^{n}\) 且 \(\varphi\) 是自同构 \((z_{1}, \ldots, z_{n}) \mapsto (z_{2}, \ldots, z_{n}, z_{1})\) 的情形；这时 \(\varphi\) 的不动点是诸点 \((z, z, \ldots, z)\)（\(z \in U\)），它们组成 T 的一个连通子群。

命题 3. 设 \(u\) 是 \(G\) 的一个自同构，并设 \(x\) 是 \(G\) 中被 \(u\) 固定的一个点。

\begin{enumerate}
\def\labelenumi{\alph{enumi})}
\item
  存在元素 \(a\)（\(\mathfrak{g}\) 中的），它被 \(\mathrm{L}(u)\) 与 \(\operatorname{Ad} x\) 固定，且使得 \(x \exp a\) 是正则的。
\item
  存在正则元 \(g\)（\(G\) 的元素），它被 \(u\) 固定且与 \(x\) 交换。
\end{enumerate}

设 H 是 u 的不动点所成的群，S 是 \(\mathrm{Z}(x)\cap\mathrm{H}\) 的一个极大环面，K 是 \(\mathrm{Z}(\mathrm{S})\cap\mathrm{Z}(x)\) 的单位元连通分支。这是 G 的一个连通闭子群；此外，由 §2 第 2 节定理 2 的推论 5，存在 G 的包含 S 与 x 的极大环面，故 K 是极大秩的且包含 S 与 x。另一方面，K 在 u 下稳定，因为 S 与 x 都在 u 下稳定；用 V 表示 u 在 K 中的不动点所成的集合。则

\[
\mathrm{S} \subset \mathrm{V} _ {0} \subset \mathrm{K} \cap \mathrm{H} \subset \mathrm{Z} (\mathrm{S}) \cap \mathrm{Z} (x) \cap \mathrm{H},
\]

于是 \(V_{0}\) 含于 S 在 \((\mathrm{Z}(x)\cap\mathrm{H})_{0}\) 中的中心化子；而后者就是 S（同前处，推论 6），故最终有 \(V_{0}=S\)。于是引理 1 蕴含 K 是交换的，故是 G 的一个极大环面（因为它连通且极大秩）。它包含 S 与 x，且等于 \(\mathrm{Z}(\mathrm{S})\cap\mathrm{Z}(x)\) 的单位元连通分支；断言 a) 便由引理 2 得出。取 \(g=x\exp a\) 即得 b)。

推论. 设 \(\mathfrak{s}\) 是一个紧李代数，并设 \(\varphi\) 是 \(\mathfrak{s}\) 的一个自同构。则存在 \(\mathfrak{s}\) 的一个被 \(\varphi\) 固定的正则元。

把 \(\mathfrak{s}\) 换成 \(\mathcal{D}\mathfrak{s}\)，可以假设 \(\mathfrak{s}\) 是半单的。设 \(S\) 是一个以 \(\mathfrak{s}\) 为李代数的紧单连通李群，并设 \(u\) 是 \(S\) 的一个自同构，满足 \(L(u) = \varphi\)。命题 3 蕴含存在元素 \(a\)（\(\mathfrak{s}\) 中的），它被 \(\varphi\) 固定，且使 \(\exp a\) 在 \(S\) 中正则；特别地，\(a\) 在 \(\mathfrak{s}\) 中正则（第 2 节，注 4）。

定理 1. 设 \(u\) 是连通紧李群 \(G\) 的一个自同构。

\begin{enumerate}
\def\labelenumi{\alph{enumi})}
\item
  \(u\) 的不动点所成的群的单位元连通分支包含 \(G\) 的一个正则元。
\item
  存在 G 的一个极大环面 K 与 L(K) 的一个房，它们在 u 下稳定。
\item
  若 G 是单连通的，则 u 的不动点所成的集合是连通的。
\end{enumerate}

断言 a) 是命题 3 当 x = e 时的特款。现在假设 G 单连通，证明 b) 与 c)。设 x 是 G 的被 u 固定的一个元素，并设 g 是 G 的一个正则元，它被 u 固定且与 x 交换（命题 3）。g 的中心化子 K 是 G 的一个极大环面（第 2 节，注 2），它在 u 下稳定，且包含 x 与 g。由第 2 节命题 2 的推论 3，存在 L(K) 的唯一单房 A 使得 \(g \in \exp A\) 且 \(0 \in \overline{A}\)；由于 g 被 u 固定，L(u) 使 A、从而也使 L(K) 中包含 A 的房保持稳定。这就证明了 b)；此外，K 中被 u 固定的点所成的集合是连通的（引理 3）且包含 x 与 e，故得 c)（《一般拓扑》，第一章，§11，第 1 节，命题 2）。

剩下在一般情形证明 \(b)\)。设 \(\tilde{\mathrm{D}}(\mathrm{G})\) 是 \(\mathrm{D}(\mathrm{G})\) 的泛覆盖，且 \(f: \tilde{\mathrm{D}}(\mathrm{G}) \to \mathrm{G}\) 是典则态射，则存在 \(\tilde{u}\)（\(\tilde{\mathrm{D}}(\mathrm{G})\) 的一个自同构）使得 \(f \circ \tilde{u} = u \circ f\)。若 \(\tilde{\mathrm{K}}\) 是 \(\tilde{\mathrm{D}}(\mathrm{G})\) 的一个极大环面，而 \(\tilde{\mathrm{C}}\) 是 \(\mathrm{L}(\tilde{\mathrm{K}})\) 的一个在 \(\tilde{u}\) 下稳定的房（由已证结果，这样的房存在），则存在（§2，第 3 节，注 2）唯一的极大环面 \(\mathrm{K}\)（\(\mathrm{G}\) 的）与唯一的房 \(\mathrm{C}\)（\(\mathrm{L}(\mathrm{K})\) 的），使得 \(\tilde{\mathrm{K}} = f^{-1}(\mathrm{K})\) 且 \(\tilde{\mathrm{C}} = \mathrm{L}(f)^{-1}(\mathrm{C})\)，我们立即看出 \(\mathrm{K}\) 与 \(\mathrm{C}\) 在 \(u\) 下稳定，故断言 \(b)\) 在一般情形成立。

推论 1. 假设 Z-模 \(\pi_{1}(G)\) 是无挠的。

\begin{enumerate}
\def\labelenumi{\alph{enumi})}
\item
  G 的每个元素的中心化子都是连通的。
\item
  G 中任意两个交换的元素属于同一个极大环面。
\end{enumerate}

由 §4 第 6 节命题 11 的推论 3，D(G) 是单连通的。我们有 G = C(G) \(_0\) .D(G)；设 \(x \in\) G；写 \(x = uv\)，其中 \(u \in\) C(G) \(_0\)，\(v \in\) D(G)。则 Z(x) = C(G) \(_0\) .Z \(_{\text{D(G)}}\) (v)。由定理 1 c)，Z \(_{\text{D(G)}}\) (v) 是连通的，故 Z(x) 连通，由此得 a)。于是，由 §2 第 2 节定理 2 的推论 3，Z(x) 是 G 的包含 \(x\) 的诸极大环面之并，由此得 b)。

推论 2. 设 \(\Gamma\) 是 \(\operatorname{Aut}(\mathbf{G})\) 的一个具有以下性质的紧子群：

(*) 存在 \(u_{1}, \ldots, u_{n}\)（\(\Gamma\) 的元素），使得对所有 \(i\)，\(\Gamma_{i}\)（\(\Gamma\) 中由 \(u_{1}, \ldots, u_{i}\) 生成的子群的闭包）是 \(\Gamma\) 的正规子群，且 \(\Gamma_{n} = \Gamma\)。

则存在 G 的一个在 \(\Gamma\) 的作用下稳定的极大环面。

我们对 G 的维数用归纳法论证。显然，可以假设 \(u_{1} \neq Id\)；则 \(u_{1}\) 的不动点所成的子群 H 异于 G，且在 \(\varGamma\) 的作用下稳定。此外，由于 \(\varGamma\) 紧，\(\varGamma\) 在 \(\mathrm{Aut}(\mathrm{H}_{0})\) 中的像是 \(\varGamma\) 的一个商，故它也满足条件 (*)。由归纳假设，存在 H 的一个在 \(\varGamma\) 下稳定的极大环面 S。S 在 G 中的中心化子 K 是连通的（\(\S2\)，第 2 节，推论 5）且在 \(\varGamma\) 下稳定；它是 G 的一个极大环面，因为 \(H_{0}\) 包含 G 的一个正则元（定理 1 a)），而该元素共轭于 S 的某个元素（同前处，推论 4）。

推论 3. 设 H 是一个李群，\(\Gamma\) 是 H 的一个紧子群。假设 \(\mathrm{H}_0\) 是紧的，且 \(\Gamma\) 满足推论 2 的条件 (*)。则存在 \(\mathrm{H}_0\) 的一个极大环面 T 使得 \(\Gamma \subset \mathrm{N}_{\mathrm{H}}(\mathrm{T})\)。

推论 4. 紧李群的每个幂零子群都含于某个极大环面的正规化子中。

设 H 是一个紧李群，N 是 H 的一个幂零子群。则 N 的闭包 \(\varGamma\) 也是一个幂零群（第三章，§9，第 1 节，命题 1 的推论 2），且鉴于推论 3，只需证明 \(\varGamma\) 满足条件 (*)。现在 \(\varGamma_{0}\) 是一个连通紧幂零李群，故是一个环面（§1，第 4 节，命题 4 的推论 1），且存在元素 \(u_{1}\)（\(\varGamma\) 的），它生成 \(\varGamma_{0}\) 的一个稠密子群（《一般拓扑》，第七章，§1，第 3 节，命题 8 之前的正文）。有限群 \(\varGamma/\varGamma_{0}\) 是幂零的，且存在 \(\tilde{u}_{2},\ldots,\tilde{u}_{n}\in\varGamma/\varGamma_{0}\) 生成 \(\varGamma/\varGamma_{0}\)，且 \(\varGamma/\varGamma_{0}\) 中由 \((\tilde{u}_{2},\ldots,\tilde{u}_{r})\) 生成的子群对 \(r=2,\ldots,n\) 是正规的（《代数》，第一章，§6，第 5 节，定理 1 与第 7 节，定理 4）。于是，若 \(u_{2},\ldots,u_{n}\) 是 \(\tilde{u}_{2},\ldots,\tilde{u}_{n}\) 在 \(\varGamma\) 中的代表，则序列 \((u_{1},\ldots,u_{n})\) 具有所要求的性质。

例. 取 \(\mathbf{G} = \mathbf{U}(n,\mathbf{C})\)。我们稍后将看到，\(\mathbf{G}\) 中由对角矩阵组成的子群是 \(\mathbf{G}\) 的一个极大环面，且其正规化子是 \(\mathbf{G}\) 中由单项矩阵组成的集合（《代数》，第二章，§10，第 7 节，例 II）。

由此我们断言：若 \(\Phi\) 是有限维复向量空间 V 上一个分离的正埃尔米特型，且 \(\Gamma\) 是 \(\mathbf{U}(\Phi)\) 的一个幂零子群，则存在 V 的一个基，使 \(\Gamma\) 的诸元素的矩阵在此基下都是单项的（「Blichtfeldt 定理」）。

\subsection{\texorpdfstring{映射 \((\mathbf{G} / \mathbf{T})\times \mathbf{T}\to \mathbf{G}\) 与 \((\mathbf{G} / \mathbf{T})\times \mathbf{A}\rightarrow \mathbf{G}_r\)}{映射 (\textbackslash mathbf\{G\} / \textbackslash mathbf\{T\})\textbackslash times \textbackslash mathbf\{T\}\textbackslash to \textbackslash mathbf\{G\} 与 (\textbackslash mathbf\{G\} / \textbackslash mathbf\{T\})\textbackslash times \textbackslash mathbf\{A\}\textbackslash rightarrow \textbackslash mathbf\{G\}\_r}}

从 G × T 到 G 的映射 \((g,t)\mapsto gtg^{-1}\) 通过过渡到商诱导出解析流形间的一个态射

\[
f: (\mathrm{G} / \mathrm{T}) \times \mathrm{T} \to \mathrm{G},
\]

它是满射的（§2，第 2 节，定理 2）。通过限制，\(f\) 诱导出一个满射态射

\[
f _ {r}: (\mathrm{G} / \mathrm{T}) \times \mathrm{T} _ {r} \rightarrow \mathrm{G} _ {r}.
\]

与 \(Id_{G/T} \times exp_{T}\) 复合，我们得到同样是满射的态射

\[
\begin{array}{c} \varphi : (\mathrm{G/T}) \times \mathfrak {t} \to \mathrm{G}, \\ \varphi_ {r}: (\mathrm{G/T}) \times \mathfrak {t} _ {r} \to \mathrm{G} _ {r}; \end{array}
\]

最后，若 A 是 t 的一个单房，则 \(\varphi_{r}\) 诱导出一个满射态射

\[
\varphi_ {\mathrm{A}}: (\mathrm{G/T}) \times \mathrm{A} \to \mathrm{G} _ {r}.
\]

我们如下定义 W 在 G/T 上的一个右作用：设 \(w \in W\) 与 \(u \in G/T\)；把 w 提升为 \(\mathrm{N}_{\mathrm{G}}(\mathrm{T})\) 的一个元素 n，把 u 提升为 G 的一个元素 g。则 gn 在 G/T 中的像不依赖于 n 与 g 的选取；我们把它记作 u.w。

对这个作用，W 自由地作用于 G/T 上：实际上，沿用前面的记号，假设 u.w = u；则 \(gn \in gT\)，故 \(n \in T\) 且 w = 1。

我们定义 W 在 \((\mathrm{G}/\mathrm{T}) \times \mathrm{T}\) 上的一个右作用如下：

\[
(u, t). w = (u. w, w ^ {- 1} (t)), \quad u \in \mathrm{G} / \mathrm{T}, t \in \mathrm{T}, w \in \mathrm{W}
\]

再定义 \(W_{a}^{\prime}\) 在 (G/T) × t 上的一个右作用如下：

\[
(u, x). \omega = (u. \overline {{\omega}}, \omega^ {- 1} (x)), \quad u \in \mathrm{G/T}, x \in \mathfrak {t}, \omega \in \mathrm{W} _ {a} ^ {\prime},
\]

其中 \(\overline{\omega}\) 是 \(\omega\) 在商 \(W_{a}^{\prime}/\Gamma(T)=W\) 中的像。

若 A 是 t 的一个单房，且 \(H_{A}\) 是 \(W_{a}^{\prime}\) 中稳定 A 的子群，则通过限制得到 \(H_{A}\) 在 \((\mathrm{G}/\mathrm{T}) \times \mathrm{A}\) 上的一个作用。

这些不同的作用与态射 \(f, \varphi\) 及 \(\varphi_{\mathrm{A}}\) 相容：对 \(u \in \mathrm{G} / \mathrm{T}, t \in \mathrm{T}, x \in \mathfrak{t}, y \in \mathrm{A}, w \in \mathrm{W}, \omega \in \mathrm{W}_a', h \in \mathrm{H}_{\mathrm{A}}\)，我们有

\[
f ((u, t). w) = f (u, t), \varphi ((u, x). \omega) = \varphi (u, x), \varphi_ {\mathrm{A}} ((u, y). h) = \varphi_ {\mathrm{A}} (u, y).
\]

引理 4. 设 \(g \in \mathrm{G}\)，\(t \in \mathrm{T}\)，并设 \(\bar{g}\) 是 \(g\) 在 \(\mathrm{G} / \mathrm{T}\) 中的像。把 \(\mathrm{G} / \mathrm{T}\)（相应地 T，相应地 G）在 \(\bar{g}\)（相应地 t，相应地 \(gtg^{-1}\)）处的切空间与 \(\mathfrak{g} / \mathfrak{t}\)（相应地 \(\mathfrak{t}\)，相应地 \(\mathfrak{g}\)）等同起来，其借助的是左平移 \(\gamma(g)\)（由 \(g\)（相应地 \(t\)，相应地 \(gtg^{-1}\)）所作）。于是，\(f\) 在 \((\bar{g}, t)\) 处的切线性映射就等同于按下述方式定义的线性映射 \(f': (\mathfrak{g} / \mathfrak{t}) \times \mathfrak{t} \to \mathfrak{g}\)：若 \(z \in \mathfrak{g}, x \in \mathfrak{t}\)，且用 \(\bar{z}\) 表示 \(z\) 在 \(\mathfrak{g} / \mathfrak{t}\) 中的像，则

\[
f ^ {\prime} (\bar {z}, x) = (\mathrm{Ad} g t ^ {- 1}) (z - (\mathrm{Ad} t) z + x).
\]

设 \(\mathrm{F}\) 是从 \(\mathrm{G} \times \mathrm{T}\) 到 \(\mathrm{T}\) 的映射，满足 \(\mathrm{F}(g,t) = gtg^{-1}\)。由于 \(\mathrm{F} \circ (\gamma(g), \mathrm{Id}_{\mathrm{T}}) = \operatorname{Int} g \circ \mathrm{F}\)，我们有 \(\mathrm{T}_{(g,t)}(\mathrm{F})(gz, tx) = \mathrm{T}_t(\operatorname{Int} g) \circ \mathrm{T}_{(e,t)}(\mathrm{F})(z, tx)\)；由第三章 §3 第 12 节命题 46，

\[
\mathrm{T} _ {(e, t)} (\mathrm{F}) (z, t x) = t ((\mathrm{Ad} t ^ {- 1}) z - z) + t x = t ((\mathrm{Ad} t ^ {- 1}) (z - (\mathrm{Ad} t) z + x))
\]

从而

\[
\mathrm{T} _ {(g, t)} (\mathrm{F}) (g z, t x) = g t g ^ {- 1} ((\mathrm{Ad}   g t ^ {- 1}) (z - (\mathrm{Ad}   t) z + x)).
\]

对这个公式过渡到商，立即得到引理。

命题 4. a) 设 \(g \in \mathbf{G}, t \in \mathrm{T}, x \in \mathfrak{t}\)，并设 \(\bar{g}\) 是 \(g\) 在 \(\mathrm{G} / \mathrm{T}\) 中的像。下列条件等价：

\begin{enumerate}
\def\labelenumi{(\roman{enumi})}
\tightlist
\item
  \(t \in \mathrm{T}_r\)（相应地 \(x \in \mathfrak{t}_r\)）。
\end{enumerate}

(i bis) 元素 \(f(\bar{g}, t)\)（相应地 \(\varphi(\bar{g}, x)\)）在 G 中是正则的。

\begin{enumerate}
\def\labelenumi{(\roman{enumi})}
\setcounter{enumi}{1}
\tightlist
\item
  映射 \(f\)（相应地 \(\varphi\)）在点 \((\bar{g}, t)\)（相应地 \((\bar{g}, x)\)）处是浸没。
\end{enumerate}

(ii bis) 映射 \(f\)（相应地 \(\varphi\)）在点 \((\bar{g}, t)\)（相应地 \((\bar{g}, x)\)）处是平展的。

\begin{enumerate}
\def\labelenumi{\alph{enumi})}
\setcounter{enumi}{1}
\item
  映射 \(f_{r}\)（相应地 \(\varphi_{r}\)，相应地 \(\varphi_{A}\)）使 \((\mathrm{G}/\mathrm{T}) \times \mathrm{T}_{r}\)（相应地 \((\mathrm{G}/\mathrm{T}) \times \mathfrak{t}_{r}\)，相应地 \((\mathrm{G}/\mathrm{T}) \times \mathrm{A}\)）成为 \(G_{r}\) 的一个以 W（相应地 \(W_{a}^{\prime}\)，相应地 \(H_{A}\)）为群的主覆盖。
\item
  (i) 与 (i bis) 的等价性是明显的；(ii) 与 (ii bis) 的等价性由关系式 \(\dim((\mathrm{G}/\mathrm{T})\times\mathrm{T})=\dim((\mathrm{G}/\mathrm{T})\times\mathfrak{t})=\dim(\mathrm{G})\) 得出。由引理 4，f 在点 \((\bar{g},t)\) 处是浸没当且仅当 \(g=t+Im(Ad t-Id)\)，而这意味着 t 是正则的。最后，由于 \(\varphi=f\circ(Id_{G/T}\times\exp_{T})\)，\(\varphi\) 在点 \((\bar{g},x)\) 处是平展的当且仅当 f 在点 \((\bar{g},\exp x)\) 处是平展的，由前所述，这就是说 x 属于 \(t_{r}\)。
\item
  于是，态射 \(f_{r}, \varphi_{r}, \varphi_{A}\) 都是平展的。另一方面，W 自由地作用于 G/T 上，从而更是自由地作用于 \((\mathrm{G}/\mathrm{T}) \times \mathrm{T}\) 上。设 \(g, g'\)（G 中）与 \(t, t'\)（\(T_{r}\) 中）满足 \(f(\bar{g}, t) = f(\bar{g}', t')\)；则 \(\operatorname{Int} g^{-1} g'\) 把 \(t'\) 映到 t，从而正规化 T，因为 \(\mathrm{T} = \mathrm{Z}(t)_{0} = \mathrm{Z}(t')_{0}\)，且 \(g^{-1} g'\) 在 W 中的类 w 把 \((\bar{g}, t)\) 映到 \((\bar{g}', t')\)。由此可知，\(f_{r}\) 是以 W 为群的主覆盖；这立即蕴含 \(\varphi_{r}\) 是以 \(W_{a}'\) 为群的主覆盖，从而通过限制到 \((\mathrm{G}/\mathrm{T}) \times \mathrm{A}\)（\((\mathrm{G}/\mathrm{T}) \times \mathbf{t}_{r}\) 的连通分支）上，得到 \(\varphi_{A}\) 是以 \(H_{A}\) 为群的主覆盖。
\end{enumerate}

注. 1) 由 §2 第 4 节命题 3，流形 \((\mathrm{G}/\mathrm{T}) \times \mathrm{A}\) 是单连通的。由此可知，\(\varphi_{A}\) 是 \(G_{r}\) 的一个泛覆盖；由于 \(\pi_{1}(G_{r})\) 典则地同构于 \(\pi_{1}(G)\)（第 1 节，命题 1 与《一般拓扑》，

第十一章，编写中），我们重新得到 \(\pi_{1}(\mathrm{G})\) 可等同于 \(H_{A}\)（即等同于 \(\Gamma(\mathrm{T})/\mathrm{N}(\mathrm{G},\mathrm{T})\)）这一事实。

\begin{enumerate}
\def\labelenumi{\arabic{enumi})}
\setcounter{enumi}{1}
\item
  \(\varphi_{A}\) 在 \(W \times A \subset (G/T) \times A\) 上的限制使 \(W \times A\) 成为 \(T_{r}\) 的一个以 \(H_{A}\) 为群的主覆盖。于是我们重新得到第 2 节命题 2 的推论 1。
\item
  用 \(\mathfrak{g}_r\) 表示在指数映射下 \(\mathrm{G}_r\) 的逆像，并用 \(\varepsilon : \mathfrak{g}_r \to \mathrm{G}_r\) 表示由 \(\exp_{\mathrm{G}}\) 诱导的映射。映射 \((g, x) \mapsto (\operatorname{Ad} g)(x)\)（从 \(\mathrm{G} \times \mathfrak{t}_r\) 到 \(\mathfrak{g}_r\)）通过过渡到商定义出一个映射 \(\psi_r : (\mathrm{G} / \mathrm{T}) \times \mathfrak{t}_r \to \mathfrak{g}_r\)。我们有 \(\varepsilon \circ \psi_r = \varphi_r\)。设 \(w \in \mathrm{W}, \gamma \in \Gamma(\mathrm{T})\) 与 \(\omega \in \mathrm{W}_a'\) 满足 \(\omega(z) = w(z) + \gamma\)（对所有 \(z \in \mathfrak{t}\)）；则 \(\psi_r((\bar{g}, x)\omega) = \psi_r(\bar{g}, x) - (\operatorname{Ad} g)(\gamma)\) 对 \(g \in \mathrm{G}, x \in \mathfrak{t}_r\) 成立，故 \(\psi_r((\bar{g}, x)\omega) = \psi_r(\bar{g}, x)\) 当且仅当 \(\gamma = 0\)。由此可知（参看《一般拓扑》，第十一章，编写中），\(\psi_r\) 是 \(\mathfrak{g}_r\) 的一个以 W 为群的主覆盖，而 \(\varepsilon : \mathfrak{g}_r \to \mathrm{G}_r\) 是相伴于主覆盖 \(\varphi_r\) 的一个覆盖，其纤维同构于 \(\mathrm{W}_a'\)-集合 \(\mathrm{W}_a'/\mathrm{W}\)。
\end{enumerate}

\section{紧李群上的积分}

我们沿用 §4 的记号；令 \(w(\mathrm{G}) = \mathrm{Card}(\mathrm{W}_{\mathrm{G}}(\mathrm{T}))\)。用 dg（相应地 dt）表示 G（相应地 T）上总质量为 1 的 Haar 测度，并用 n（相应地 r）表示 G（相应地 T）的维数。

\subsection{交错多重线性形式的乘积}

设 A 是一个交换环，M 是一个 A-模。对每个整数 \(r \geq 0\)，用 \(\operatorname{Alt}^{r}(\mathbf{M})\) 表示 M 上交错 r-线性形式所成的 A-模；它可等同于 A-模 \(\Lambda^{r}(\mathbf{M})\) 的对偶（《代数》，第三章，§7，第 4 节，命题 7）。设 \(u \in \operatorname{Alt}^{s}(\mathbf{M})\) 与 \(v \in \operatorname{Alt}^{r}(\mathbf{M})\)；回忆（《代数》，第三章，§11，第 2 节，例 3）u 与 v 的交错积是由下式定义的元素 \(u \wedge v \in \operatorname{Alt}^{s+r}(\mathbf{M})\)：

\[
(u \wedge v) (x _ {1}, \dots , x _ {s + r}) = \sum_ {\sigma \in \mathfrak {S} _ {s, r}} \varepsilon_ {\sigma} u (x _ {\sigma (1)}, \dots , x _ {\sigma (s)}) v (x _ {\sigma (s + 1)}, \dots , x _ {\sigma (s + r)}),
\]

其中 \(S_{s,r}\) 是对称群 \(S_{s+r}\) 中由在 \((1,s)\) 与 \((s+1,s+r)\) 上的限制都是递增的置换所组成的子集。

现在设

\[
0 \longrightarrow \mathrm{M} ^ {\prime} \stackrel {i} {\longrightarrow} \mathrm{M} \stackrel {p} {\longrightarrow} \mathrm{M} ^ {\prime \prime} \longrightarrow 0
\]

是自由 A-模的一个正合序列，其秩分别为 \(r, r + s\) 与 \(s\)。

引理 1. 存在一个从 \(\mathrm{Alt}^s (\mathbf{M}'')\times \mathrm{Alt}^r (\mathbf{M}')\) 到 \(\mathrm{Alt}^{s + r}(\mathbf{M})\) 的 A-双线性映射，记作 \((u,v)\mapsto u\cap v\)，它由下列两个性质之一刻画：

\begin{enumerate}
\def\labelenumi{\alph{enumi})}
\item
  用 \(u_{1} \in \operatorname{Alt}^{s}(\operatorname{M})\) 表示形式 \((x_{1}, \ldots, x_{s}) \mapsto u(p(x_{1}), \ldots, p(x_{s}))\)，并设 \(v_{1} \in \operatorname{Alt}^{r}(\operatorname{M})\) 是一个满足 \(v_{1}(i(x_{1}^{\prime}), \ldots, i(x_{r}^{\prime})) = v(x_{1}^{\prime}, \ldots, x_{r}^{\prime})\)（对 \(x_{1}^{\prime}, \ldots, x_{r}^{\prime}\)，\(M^{\prime}\) 中的）的形式；则 \(u \cap v = u_{1} \wedge v_{1}\)。
\item
  对所有 \(x_{1},\ldots,x_{s}\)（M 中）与 \(x_{1}^{\prime},\ldots,x_{r}^{\prime}\)（\(M^{\prime}\) 中），
\end{enumerate}

\[
(u \cap v) (x _ {1}, \dots , x _ {s}, i (x _ {1} ^ {\prime}), \dots , i (x _ {r} ^ {\prime})) = u (p (x _ {1}), \dots , p (x _ {s})) v (x _ {1} ^ {\prime}, \dots , x _ {r} ^ {\prime}).\tag{1}
\]

映射 \(\varphi : \mathrm{Alt}^s (\mathrm{M}'')\otimes_{\mathrm{A}}\mathrm{Alt}^r (\mathrm{M}')\to \mathrm{Alt}^{s + r}(\mathrm{M})\)（满足 \(\varphi (u\otimes v) = u\cap v\)）是秩 1 自由 A-模之间的一个同构。

形式 \(v_{1}\)（满足条件 \(a\))）的存在性来自下述事实：\(\Lambda^r (i)\) 诱导出从 \(\Lambda^r (\mathrm{M}')\) 到 \(\Lambda^r (\mathrm{M})\) 的一个直因子子模的同构（《代数》，第三章，§7，第 2 节）。设 \(v_{1}\) 是这样一个形式；令 \(u\cap v = u_{1}\wedge v_{1}\)。这时公式 (1) 得到满足，因为若令 \(i(x_k') = x_{s + k}\)（对 \(1\leq k\leq r\)），则元素 \(\sigma\)（属于 \(\mathfrak{S}_{s,r}\)）中使 \(p(x_{\sigma (i)})\neq 0\)（对 \(1\leq i\leq s\)）的只有恒等置换。另一方面，公式 (1) 唯一确定 \(u\cap v\)：实际上，设 \((e_1',\ldots ,e_r')\) 是 \(\mathrm{M}'\) 的一个基，\((f_1'',\dots,f_s'')\) 是 \(\mathrm{M}''\) 的一个基，而 \(f_{1},\ldots ,f_{s}\) 是 \(\mathrm{M}\) 中满足 \(p(f_i) = f_i''\)（对 \(1\leq i\leq s\)）的元素。则 \((f_{1},\ldots ,f_{s},i(e_1'),\ldots ,i(e_r'))\) 是 \(\mathrm{M}\) 的一个基（《代数》，第二章，§1，第 11 节，命题 21），而公式 (1) 可写成

\[
(u \cap v) (f _ {1}, \dots , f _ {s}, i (e _ {1} ^ {\prime}), \dots , i (e _ {r} ^ {\prime})) = u (f _ {1} ^ {\prime \prime}, \dots , f _ {s} ^ {\prime \prime}) v (e _ {1} ^ {\prime}, \dots , e _ {r} ^ {\prime});\tag{2}
\]

但 \(\mathrm{Alt}^{s + r}(\mathbf{M})\) 的一个元素由它在一个基上的值唯一确定。

由前述可知，条件 a) 与 b) 中的每一个都唯一确定乘积 \(u \cap v\)；显然这个乘积是双线性的。最后，引理的最后一个断言由公式 (2) 得出。